\documentclass[11pt,letterpaper]{article}
\usepackage[T1]{fontenc}
\usepackage{lmodern}
\usepackage[margin=1in]{geometry}
\usepackage{amsmath,amssymb,amsthm,mathtools}
\usepackage{booktabs,longtable,array,tabularx}
\usepackage{graphicx,xcolor}
\usepackage{microtype}
\usepackage{enumitem}
\usepackage{xurl}
\usepackage{hyperref}
\usepackage{fancyhdr}
\definecolor{linkblue}{RGB}{31,67,101}
\hypersetup{colorlinks=true,linkcolor=linkblue,citecolor=linkblue,urlcolor=linkblue,
 pdftitle={Three advances on OEIS conjectures and asymptotics},
 pdfauthor={Prepared with ChatGPT},
 pdfsubject={A343093; A088714 and A088713; A301981 and A301982}}
\setlength{\emergencystretch}{2em}
\setlength{\headheight}{14pt}
\setlist{topsep=4pt,itemsep=3pt,parsep=0pt}
\allowdisplaybreaks[2]
\numberwithin{equation}{section}
\newtheorem{theorem}{Theorem}[section]
\newtheorem{lemma}[theorem]{Lemma}
\newtheorem{proposition}[theorem]{Proposition}
\newtheorem{corollary}[theorem]{Corollary}
\theoremstyle{definition}
\newtheorem{definition}[theorem]{Definition}
\newtheorem{question}[theorem]{Research question}
\theoremstyle{remark}
\newtheorem{remark}[theorem]{Remark}
\newcommand{\R}{\mathbb R}
\newcommand{\C}{\mathbb C}
\newcommand{\N}{\mathbb N}
\newcommand{\eps}{\varepsilon}
\newcommand{\dd}{\,\mathrm d}
\DeclareMathOperator{\Rea}{Re}
\DeclareMathOperator{\Ima}{Im}
\DeclareMathOperator{\supp}{supp}
\DeclareMathOperator{\Res}{Res}
\pagestyle{fancy}
\fancyhf{}
\fancyhead[L]{\small Three advances on OEIS conjectures}
\fancyhead[R]{\small Research article}
\fancyfoot[C]{\thepage}
\renewcommand{\headrulewidth}{0.3pt}
\setcounter{tocdepth}{1}
\title{\textbf{Three advances on OEIS conjectures\\and asymptotics}\\[7pt]
 \large Bridgeless toroidal maps, golden-ratio moment laws,\\
 and unitary-divisor partitions}
\author{Prepared with ChatGPT\\[2pt]
 \small Building on Vladimir Reshetnikov's ProveIt research repository}
\date{4 October 2026 (UTC)}
\begin{document}
\maketitle
\begin{abstract}
We prove the square-convolution generating function posted as a conjecture
for the bridgeless toroidal maps of OEIS A343093, and derive positive
coefficient formulas and a three-term asymptotic expansion. For the
nonlinear reversion pair A088714 and A088713, building on the ProveIt
moment-law construction, we prove absolute continuity, strictly positive
real-analytic densities on the full positive half-line, and sharp density
and derivative equivalents at the origin. An inverse Fibonacci-type
iteration produces a nonzero real-analytic antiperiodic profile governing
the first correction to the golden-ratio transform law, together with
a convergent inverse-logarithmic expansion and controlled exponentially
smaller errors. A general renewal
theorem supplies all fixed-shift corrections in the comparison of the
two coefficient sequences. Finally, we strengthen the existing refutation
of the pure asymptotic models for A301981 and A301982: each coefficient-to-model
ratio has lower limit zero and upper limit infinity, with two-sided
logarithmic oscillations of at least order $n^{1/4}$ and corresponding
inverse errors. We state the prior dependencies, give complete arguments
for the advances, include reproducible exact and interval checks, and
identify further research directions.
\end{abstract}
\clearpage
\tableofcontents
\clearpage

% Included source: sections/introduction.tex
\section{Results, provenance, and the proof boundary}
\label{intro:results}

This article develops three advances on concrete OEIS questions. The first
proves a generating-function identity explicitly marked conjectural in
A343093. The second answers three analytic and asymptotic questions about
A088714 and A088713 in the ProveIt research corpus. The third strengthens
an existing refutation of two partition asymptotics to quantitative
oscillation in both directions. The subjects differ, but the proof strategy
has a common feature: an exact structural identity is retained long enough
to control information that a leading estimate loses.

The repository snapshot used throughout is
\begin{center}\small
 \texttt{eaf08931cbfd0132dd36ab82a7743ccb02506dc7},
 \quad 4 October 2026, 00:59:12 UTC.
\end{center}
The related OEIS entries and the cited primary literature were inspected
on 4 October 2026 UTC. The manuscript supplies written proofs, independent
internal mathematical reviews, exact finite checks, and an interval
certificate. These are distinct forms of evidence: a finite computation
does not establish an infinite asymptotic statement. External peer review
and proof-assistant formalization remain future work.

\subsection{A map identity, with all root multiplicities included}

If $b_n$ counts rooted toroidal maps with $n$ edges and no bridges, then
\begin{equation}\label{intro:maps}
 \sum_{n\ge2}b_nz^n
 =\left(\frac14\sum_{n\ge1}\binom{4n}{n}z^n\right)^2.
\end{equation}
This is Peter Bala's posted conjecture \cite{map:oeis343093}.
The new derivation applies the bridge decomposition to the classical
planar and toroidal generating functions of Bender--Canfield--Robinson.
The genus carried by a branch changes the substituted argument, producing
an essential derivative term. Keeping that term gives the exact square.
The resulting coefficient formula also yields
\[
 b_n=\left(\frac{256}{27}\right)^n
 \left\{\frac1{24}-\frac{7\sqrt6}{216\sqrt{\pi n}}
              +O(n^{-3/2})\right\}.
\]
The article includes the next explicit correction, an algebraic equation,
and an independent enumeration of rotation systems at small sizes.

\subsection{Three answers for the golden-ratio moment pair}

Let $A(z)=1+zA(z)^2A(zA(z))$ and $C(z)=(1-zA(z))^{-1}$.
The existing moment-law construction \cite{goldenprior} associates
probability measures $\mu,\nu$ with their coefficients. We prove that
both measures have strictly positive real-analytic densities throughout
$(0,\infty)$ and support exactly $[0,\infty)$. With
$\beta=(\sqrt5-1)/2$ and $\alpha=1-\beta$, their densities satisfy
\[
 p(x)\sim\frac{\sin(\pi\alpha)}{\pi}x^{-\beta},\qquad
 q(x)\sim\frac{\sin(\pi\beta)}{\pi}x^{-\alpha}.
\]
The proof uses a coupled inverse relation to rule out every hypothetical
positive density zero. A complex inversion argument establishes the
endpoint equivalents and their fixed derivatives.

The second answer reveals information hidden by the leading equivalent
$F(s)\sim s^{-\alpha}$ for
$F(s)=\int(1+st)^{-1}\,d\mu(t)$. There is a nonzero real-analytic function
$P$ with $P(t+1)=-P(t)$ such that
\begin{equation}\label{intro:oscillation}
 \log\bigl(s^\alpha F(s)\bigr)
 =\frac{P(\log_\varphi\log s)}{\log s}
       +O((\log s)^{-3}),\qquad \varphi=\frac{1+\sqrt5}{2}.
\end{equation}
An exact inverse-iteration construction defines the profile; an interval
calculation proves that it is not zero. An explicit expansion to all
inverse-logarithmic orders converges for sufficiently large $\log s$,
with an exponentially smaller discrepancy that is separately bounded.
The sign reversal is a theorem about the profile, rather than an
interpretation of a numerical plot.

The third answer concerns the coefficients directly. If
$d_j=[z^j]C(z)^2$, then for every fixed $K$,
\begin{equation}\label{intro:renewal}
 c_n=\sum_{j=0}^{K}d_ja_{n-1-j}+O_K(a_{n-K-2}).
\end{equation}
In particular $c_n/a_{n-1}\to1$, with first relative correction
$2\log n/n$. A general weighted-composition argument proves this without
assuming a fine expansion of $a_n/a_{n-1}$.

\subsection{Partition asymptotics with arbitrarily large errors of both signs}

Let $a_n^-$ and $a_n^+$ be the coefficients of
\[
 \prod_{j\ge1}(1-z^j)^{-\sigma^*(j)},\qquad
 \prod_{j\ge1}(1+z^j)^{\sigma^*(j)},
\]
respectively. Their pure OEIS models $m_n^\pm$ are defined precisely in
Part III. The prior report \cite{unitaryprior} already proves that these
models are not asymptotic equivalents. The new result is
\begin{equation}\label{intro:partitions}
 \liminf\frac{a_n^\pm}{m_n^\pm}=0,
 \qquad \limsup\frac{a_n^\pm}{m_n^\pm}=\infty,
 \qquad
 \log(a_n^\pm/m_n^\pm)=\Omega_\pm(n^{1/4}).
\end{equation}
The notation $\Omega_\pm$ means that positive and negative multiples of
the indicated scale occur along separate infinite subsequences.
The proof gives sharper lower constants from one noncancelled Mellin pole.
It requires neither the Riemann hypothesis nor simplicity of zeta zeros.
The error of inverting either pure model likewise has both signs on a
scale at least $(\log y)^{7/8}$.

\subsection{What the novelty claim means}

The map assertion is a proof of a currently posted conjectural identity;
the underlying map formulas and bridge decomposition are credited
classical ingredients. The moment-law existence theorem, leading
golden-ratio equivalents, and coefficient ratio law are prior ProveIt
results. The new statements resolve that report's Questions 23.1, 23.3,
and 23.6. The two-sided partition theorem answers its predecessor's
Section 11, question 3; it does not claim the earlier refutation as new.
The general methods of analytic inversion, Lagrange inversion, and
Landau positivity are also established methods.

The search did not identify earlier proofs of these particular new
conclusions. It does not establish exhaustive historical priority, and
older map-enumeration transformations may encode equivalent identities.
Each Part states its dependencies explicitly so that its correctness can
be assessed independently of any claim about importance or novelty.


\clearpage
\part{Bridgeless toroidal maps}

% Included source: sections/maps.tex
\section{A square convolution for bridgeless toroidal maps}
\label{map:section}

Let $b_n$ be the number of rooted toroidal maps with $n$ edges and no
isthmuses.  These are the numbers in OEIS A343093, with $b_0=b_1=0$.
The entry attributes the following generating-function conjecture to Peter
Bala, dated July 24, 2025 \cite{map:oeis343093}:
\begin{equation}
 \label{map:bala}
 \sum_{n\geq 2}b_n t^n
 =\left(\frac14\sum_{n\geq1}\binom{4n}{n}t^n\right)^2.
\end{equation}
We prove this identity from the classical generating functions for all
planar and toroidal maps.  The proof also gives a positive finite formula,
an algebraic equation, and an expansion beyond the leading asymptotic
term.  The result is a resolution of the explicitly stated OEIS conjecture;
the underlying enumeration of unrestricted toroidal maps and the
decomposition along bridges are established ingredients.

\subsection{Conventions and the classical inputs}
\label{map:conventions}

A map is a finite connected graph, with loops and parallel edges permitted,
cellularly embedded in a closed oriented surface.  Cellular means that
every component of the complement is an open disk.  Isomorphisms preserve
the orientation.  A root is a distinguished directed edge, or dart.  The
one-vertex edgeless planar map is included by convention and has one root
corner.  An isthmus, also called a bridge, is an edge whose deletion
disconnects the underlying graph.  A loop is consequently never a bridge.
In particular, the restriction here is on edge connectivity: cut vertices
and separating cycles are allowed.

Write $m_{g,n}$ for the number of rooted maps of genus $g$ with $n$ edges,
and put
\[
 M_g(z)=\sum_{n\geq0}m_{g,n}z^n,
 \qquad M(z,u)=\sum_{g\geq0}u^gM_g(z).
\]
For a fixed $n$, only finitely many genera occur, so that these operations
can be performed in formal power series without an analytic convergence
assumption.  Let $B_g(t)$ count the nonempty bridgeless rooted maps of genus
$g$, and let $B(t,u)=\sum_{g\geq0}u^gB_g(t)$.  In particular,
$B_1(t)=\sum_{n\geq2}b_n t^n$.

The two external enumerative inputs are
\begin{equation}
 \label{map:classical}
 M_0(z)=\frac{4(2R+1)}{3(R+1)^2},\qquad
 M_1(z)=\frac{(R-1)^2}{12R^2(R+2)},\qquad
 R=\sqrt{1-12z},\quad R(0)=1.
\end{equation}
They are equations (4.20) and (4.22), respectively, on printed page 268
of Bender, Canfield, and Robinson \cite{map:bcr1988}.  Their convention of
distinguishing an edge, its direction, and its side up to arbitrary
surface homeomorphism is equivalent, for orientable surfaces, to choosing
an orientation and a dart up to orientation-preserving homeomorphism.
Thus no factor of two is needed in \eqref{map:classical}.

\subsection{Removing bridges, with the root accounted for}
\label{map:bridges}

\begin{proposition}[Bridge decomposition]
 \label{map:decomposition}
 With the preceding conventions,
 \begin{equation}
  \label{map:master}
  M(z,u)=1+zM(z,u)^2+
  B\!\left(\frac{z}{(1-zM(z,u))^2},u\right).
 \end{equation}
\end{proposition}

\begin{proof}
The edgeless map contributes $1$.  If the root edge is a bridge, cut its
edge ribbon in a regular neighborhood of the embedded graph and cap the
boundary circles of the resulting two ribbon submaps.  The components are rooted
at the corners determined by the former root edge, and their order is
determined by its direction.  A rooted corner in a nonempty oriented map
is equivalent to a rooted dart.  Either component may be the edgeless
planar map.  Conversely, any ordered pair of rooted maps can be joined
at their root corners by a new directed bridge.  This is an inverse
construction and contributes $zM^2$.

Suppose instead that the root is not a bridge.  Remove all bridges and
retain the connected component containing the root.  It is a nonempty
bridgeless rooted core, say with $k$ edges, and it has exactly $2k$ corners.
Every component discarded from the core is attached to it through a
bridge, possibly with further bridges farther from the core.  At each
corner, the cyclic order of darts at the incident vertex gives a linearly
ordered sequence of branches.  A branch consists of its attaching bridge
and the rooted map on its far side.  Its generating function is $zM$,
so one corner contributes
\[
 1+zM+(zM)^2+\cdots=\frac{1}{1-zM}.
\]
Every one of the $2k$ corners, including the corner adjoining the root,
admits this construction.  The root remains a dart of the core; there is
no further choice of a corner or a new root.  A core of size $k$ therefore
has weight $z^k(1-zM)^{-2k}$.

The decomposition and reconstruction are unique.  Indeed, graph bridges
are intrinsic, the component containing the nonbridge root is intrinsic,
and the rotations at core vertices recover both the relevant corners and
the order of branches.  Conversely, grafting ordered branches into these
corners creates precisely the bridges that are removed by the procedure.

It remains to check the genus weight.  Joining two cellular maps of genera
$g_1,g_2$ by a bridge produces a map with
\[
 V=V_1+V_2,\qquad E=E_1+E_2+1,\qquad F=F_1+F_2-1.
\]
Euler's formula gives $2-2g=2-2(g_1+g_2)$.  Thus genera add under every
bridge attachment and also in the root-bridge case.  Multiplying the
weights $u^{g_i}$ is correct.  Summing the core contribution proves
\eqref{map:master}.
\end{proof}

This is the usual canonical decomposition into a bridgeless root component
and branches.  Courtiel, Yeats, and Zeilberger give a closely related
formulation in Proposition 23 \cite{map:cyz2019}.  Their auxiliary dangling
root is included in the size, whereas our root is one of the ordinary
edges.  The proof above fixes the normalization and explains the exponent
$2k$ in \eqref{map:master} directly.

There is also a useful inverse form.  If $P=1+B(t,u)$, then
\begin{equation}
 \label{map:inverse}
 M=P(1+tP),\qquad z=\frac{t}{(1+tP)^2}.
\end{equation}
To see this, set $Q=1-zM$ in \eqref{map:master}.  It gives $P=MQ$ and
$t=z/Q^2$, whence $Q=1/(1+tP)$.  These identities are formal: the relevant
series have constant term $1$, and $t=z+O(z^2)$ is invertible for
composition.

\subsection{Extracting genus one}
\label{map:extract}

Introduce the auxiliary variable
\begin{equation}
 \label{map:parameter}
 U=\frac{1-R}{3(1+R)},\qquad
 z=\frac{U}{(1+3U)^2}.
\end{equation}
The planar and toroidal formulas \eqref{map:classical} become
\begin{equation}
 \label{map:inputs-u}
 M_0=(1-U)(1+3U),\qquad
 M_1=\frac{U^2(1+3U)}{(1+U)(1-3U)^2}.
\end{equation}
For the moment $U$ is viewed as a formal series in $z$.  At genus zero set
\[
 Q_0=1-zM_0=\frac{(1+U)^2}{1+3U},\qquad
 t_0=\frac{z}{Q_0^2}=\frac{U}{(1+U)^4}.
\]
The genus-zero coefficient of \eqref{map:master} then gives
\begin{equation}
 \label{map:planar}
 B_0(t_0)=M_0-1-zM_0^2=U-U^2-U^3.
\end{equation}
Since
\[
 \frac{d}{dU}(U-U^2-U^3)=(1+U)(1-3U),\qquad
 \frac{dt_0}{dU}=\frac{1-3U}{(1+U)^5},
\]
we obtain the particularly simple identity
\begin{equation}
 \label{map:b0-derivative}
 B_0'(t_0)=(1+U)^6.
\end{equation}

Now expand \eqref{map:master} to first order in $u$.  The coefficient of
$u$ in its substituted argument is
\[
 [u]\frac{z}{(1-zM(z,u))^2}=\frac{2z^2M_1}{Q_0^3}.
\]
Consequently
\begin{equation}
 \label{map:linear-genus}
 B_1(t_0)=M_1\left(1-2zM_0-
          \frac{2z^2B_0'(t_0)}{Q_0^3}\right).
\end{equation}
There are no other genus-one terms: either the core has genus one and all
branches are planar, or the core is planar and exactly one branch has
genus one.  Substituting the explicit formulas into the parenthesis gives
\[
 1-\frac{2U(1-U)}{1+3U}-\frac{2U^2}{1+3U}
 =\frac{1+U}{1+3U}.
\]
Together with \eqref{map:inputs-u}, this proves the following theorem.

\begin{theorem}[Exact toroidal generating function]
 \label{map:gf-theorem}
 Let $U(t)$ be the unique formal series with $U(0)=0$ satisfying
 \begin{equation}
  \label{map:quartic-tree}
  U=t(1+U)^4.
 \end{equation}
 Then
 \begin{equation}
  \label{map:gf-rational}
  B_1(t)=\frac{U(t)^2}{(1-3U(t))^2}.
 \end{equation}
 In particular, Bala's identity \eqref{map:bala} holds for every
 coefficient.
\end{theorem}

\begin{proof}
Equation \eqref{map:linear-genus} proves \eqref{map:gf-rational}, because
$t_0=z+O(z^2)$ may be used as a formal variable.  For the claimed square,
differentiate \eqref{map:quartic-tree} to obtain
\[
 \frac{U}{1-3U}=t\frac{d}{dt}\log(1+U).
\]
Lagrange inversion gives, for $n\geq1$,
\[
 [t^n]\log(1+U(t))
 =\frac1n[u^{n-1}](1+u)^{4n-1}
 =\frac{1}{4n}\binom{4n}{n}.
\]
It follows that $U/(1-3U)=\frac14\sum_{n\geq1}\binom{4n}{n}t^n$.
Squaring establishes \eqref{map:bala}.
\end{proof}

\subsection{Exact coefficients and an algebraic certificate}
\label{map:coefficients}

\begin{corollary}
 \label{map:finite-formulas}
 For $n\geq2$,
 \begin{align}
  b_n&=\frac1{16}\sum_{k=1}^{n-1}
        \binom{4k}{k}\binom{4(n-k)}{n-k},
       \label{map:convolution}\\
  b_n&=\frac1n\sum_{j=0}^{n-2}(j+1)(j+2)3^j
        \binom{4n}{n-2-j}.
       \label{map:single-sum}
 \end{align}
\end{corollary}

\begin{proof}
The first formula is the Cauchy product in \eqref{map:bala}.  For the
second, apply Lagrange inversion to $U^2/(1-3U)^2$; its derivative with
respect to $U$ is $2U/(1-3U)^3$.  Thus
\[
 b_n=\frac2n[u^{n-2}]\frac{(1+u)^{4n}}{(1-3u)^3}.
\]
Expanding the last denominator gives \eqref{map:single-sum}.
\end{proof}

The first ten values, starting at $n=2$, are
\[
 1,\ 14,\ 159,\ 1680,\ 17147,\ 171612,\ 1696491,\
 16631840,\ 162090756,\ 1572801142.
\]
The first factor in the convolution is integral:
$\frac14\binom{4k}{k}=\binom{4k-1}{k-1}$.
Hence \eqref{map:convolution} also gives an elementary positive model for
the numbers: ordered pairs of binary words of respective lengths
$4k-1$ and $4(n-k)-1$, containing $k-1$ and $n-k-1$ ones.
The enumeration proves this equinumerosity; constructing a direct
geometric bijection is a further question.

For an explicit algebraic certificate put $A=B_1(t)$ and
$H=U/(1-3U)$.  Solving for $U$ yields
\[
 t(1+4H)^4=H(1+3H)^3,\qquad A=H^2.
\]
Separating even and odd powers of $H$ gives
\begin{equation}
 \label{map:algebraic}
 \left[t+(96t-9)A+(256t-27)A^2\right]^2
 -A\left[(16t-1)+(256t-27)A\right]^2=0.
\end{equation}
The intended solution is the one specified by \eqref{map:quartic-tree}
and \eqref{map:gf-rational}.  Equation \eqref{map:algebraic} alone is an
elimination relation and should not be used to choose a branch without
these initial conditions.

\subsection{Growth and its first corrections}
\label{map:asymptotics}

\begin{theorem}[Asymptotic expansion]
 \label{map:asymptotic-theorem}
 With $\rho=27/256$, as $n\to\infty$,
 \begin{equation}
  \label{map:asymptotic}
  b_n=\rho^{-n}\left[
    \frac1{24}
    -\frac{7\sqrt6}{216\sqrt\pi}\,n^{-1/2}
    +\frac{217\sqrt6}{93312\sqrt\pi}\,n^{-3/2}
    +O(n^{-5/2})\right].
 \end{equation}
 A full expansion in the successive powers $n^{-1/2},n^{-3/2},\ldots$
 follows by the coefficient procedure in the proof.  In particular,
 $b_n\sim (256/27)^n/24$.
\end{theorem}

\begin{proof}
The characteristic equation for $U=t(1+U)^4$ is $1=4t(1+U)^3$.
It has the positive solution $U=1/3$, $t=\rho$.  Equivalently,
$t=U/(1+U)^4$ has its only finite critical value away from $t=0$ at
$\rho$.  The branch $U(0)=0$ is analytic at zero and admits analytic
continuation to a disk slightly larger than $|t|<\rho$, slit along the
positive real ray starting at $\rho$.  One can see this directly from
the algebraic equation: multiple finite roots require $U=1/3$ and
$t=\rho$, and the leading coefficient can vanish only at $t=0$,
where this particular branch is already regular.  Thus the singularity
at $\rho$ is the unique dominant one for this branch.

For explicit calculations it is convenient to reuse
$R=(1-3U)/(1+3U)$.  Then
\[
 \frac{t}{\rho}=
 \frac{16(1-R)(1+R)^3}{(R+2)^4},\qquad
 B_1(t)=\frac{(1-R)^2}{36R^2}.
\]
Set $w=\sqrt{1-t/\rho}$, positive for $0<t<\rho$.  The first relation
has a convergent local inverse of the form
\[
 R=\frac{\sqrt6}{3}w+\frac19w^2
       +\frac{19\sqrt6}{648}w^3-\frac{11}{972}w^4
       +\frac{677\sqrt6}{93312}w^5+O(w^6).
\]
The existence of this inverse follows by writing
$1-t/\rho=R^2(3/2+O(R))$ and applying the analytic inverse-function
theorem to $R\sqrt{3/2+O(R)}$.
Substitution gives
\begin{equation}
 \label{map:puiseux}
 B_1(t)=\frac{1}{24w^2}-\frac{7\sqrt6}{216w}
        +\frac{83}{2592}
        +\frac{161\sqrt6}{46656}w
        -\frac{1249}{279936}w^2+O(w^3).
\end{equation}
For nonintegral $\alpha$ the exact coefficient identity
\[
 [t^n](1-t/\rho)^\alpha
 =\rho^{-n}\frac{\Gamma(n-\alpha)}{
                    \Gamma(-\alpha)\Gamma(n+1)}
\]
and the gamma-ratio expansion convert the nonanalytic terms in
\eqref{map:puiseux} into coefficient asymptotics.  Local analytic terms
contribute no algebraic-order term from the dominant singularity.
In particular,
\begin{align*}
 [t^n](1-t/\rho)^{-1/2}
 &=\frac{\rho^{-n}}{\sqrt{\pi n}}
       \left(1-\frac1{8n}+O(n^{-2})\right),\\
 [t^n](1-t/\rho)^{1/2}
 &=-\frac{\rho^{-n}}{2\sqrt\pi\,n^{3/2}}
       \left(1+O(n^{-1})\right).
\end{align*}
The $w^{-2}$ term has coefficient exactly $\rho^{-n}/24$.
The coefficient of $\rho^{-n}n^{-3/2}$ is
\[
 \frac{7\sqrt6}{1728\sqrt\pi}
 -\frac{161\sqrt6}{93312\sqrt\pi}
 =\frac{217\sqrt6}{93312\sqrt\pi}.
\]
This proves \eqref{map:asymptotic}.  Continuing the convergent Puiseux
inversion, followed by the gamma-ratio expansion for each odd power of
$w$, gives any prescribed number of terms, with a remainder of the
corresponding order.  The slit-domain continuation supplies the usual
coefficient-transfer justification.
\end{proof}

For comparison, all toroidal maps satisfy $m_{1,n}\sim12^n/24$, as follows
from \eqref{map:classical}.  Hence the probability that a uniformly
selected rooted toroidal map with $n$ edges has no bridges satisfies
\begin{equation}
 \label{map:probability}
 \frac{b_n}{m_{1,n}}\sim\left(\frac{64}{81}\right)^n.
\end{equation}
This describes the exponential effect of the connectivity restriction.

\subsection{Verification and further questions}
\label{map:verification}

The accompanying program \texttt{code/verify\_maps.py} checks the
convolution and Lagrange formulas against all twenty terms printed in the
inspected OEIS entry, and against each other for additional indices.
It also contains an independent enumeration for small maps.  Fix the edge
involution on $2n$ darts to be
$(0\ 1)(2\ 3)\cdots(2n-2\ 2n-1)$ and enumerate every vertex permutation
$\sigma$.  Connectivity is tested in the underlying multigraph; the
cycles of $\sigma$ and $\sigma\alpha$ give vertices and faces.  Euler's
formula determines the genus, and deleting each edge tests the bridge
restriction.  Keeping dart $0$ as root, each isomorphism class is counted
exactly $2^{n-1}(n-1)!$ times: these are precisely the relabelings
commuting with the fixed involution and fixing the root.  Their action
on connected rooted maps is free.  Thus division by this known factor
gives a genuinely separate check of the normalization and the first
coefficients.  These finite checks support the implementation; the
all-index identity is established by Theorem~\ref{map:gf-theorem}.

Several questions now have concrete starting points.

\begin{enumerate}
 \item Find a direct bijection underlying the square in
 \eqref{map:bala}.  The product describes two independently chosen
 objects of positive sizes, and a geometric interpretation should explain
 how a toroidal bridgeless map yields that ordered pair.
 \item Refine the square identity by vertices and faces, using the
 two-variable map enumeration and OEIS A343092.  A valid refinement must
 reduce to \eqref{map:bala} after summing the face parameter; an arbitrary
 weighting of the two factors need not count maps.
 \item Carry \eqref{map:master} to genus two and higher.  At each genus
 the new $B_g$ enters linearly after all smaller genera are known, so the
 transformation gives an effective route from unrestricted-map formulas
 to bridgeless-map formulas.  Determining which genera admit comparably
 simple positive factorizations is a separate question.
 \item Mark bridges with a variable $q$.  The same decomposition gives
 $M=1+qzM^2+B(z/(1-qzM)^2,u)$, with $M$ now recording the bridge count.
 This supplies exact generating functions from which to study the
 distribution of the number of bridges and of the core sizes.
\end{enumerate}

The bounded source audit found the conjecture still explicitly marked
as such in the inspected OEIS entry and found no A343093 match in the
inspected ProveIt snapshot.  It also identified the classical sources
for both crucial inputs.  This supports the precise description
``a proof of the posted OEIS identity''; it does not establish that every
equivalent specialization is absent from the older map-enumeration
literature.


\clearpage
\part{The nonlinear reversion pair A088714 and A088713}

% Included source: sections/golden_setting.tex
\section{The nonlinear reversion pair and the inherited foundations}
\label{gold:setting}

The ordinary generating functions in this Part are formal series
\begin{equation}\label{gold:formal}
 A(z)=\sum_{n\ge0}a_nz^n=1+zA(z)^2A(zA(z)),
 \qquad C(z)=\sum_{n\ge0}c_nz^n=\frac1{1-zA(z)}.
\end{equation}
Here $a_n$ is A088714 and $c_n$ is A088713. Their initial values are
\[
 (a_0,\ldots,a_6)=(1,1,3,13,69,419,2809),\qquad
 (c_0,\ldots,c_6)=(1,1,2,6,24,118,674).
\]
The definitions, reversion relation, and identity $A(z)=C(zA(z))$
are recorded in the OEIS entries \cite{oeisA088714,oeisA088713}.
The factor $z$ makes the first equation triangular in degree and gives
a unique formal solution. The series have radius of convergence zero;
analytic statements below concern their Stieltjes transforms, not evaluations
of these divergent ordinary series at nonzero arguments.

The ProveIt report \cite{goldenprior}, Parts I--III, is the starting point.
Its Section 23 explicitly asks about regularity and support (Question 23.1),
second-order golden-ratio asymptotics (Question 23.3), and the sharper
companion comparison (Question 23.6). We resolve these questions below.
They are questions formulated in that research report; the live OEIS entries
themselves do not state all three as conjectures.

\begin{proposition}[Foundation imported from the pinned report]
\label{gold:foundation}
There are unique probability measures $\mu,\nu$ on $[0,\infty)$ with
\[
 a_n=\int_0^\infty t^n\,d\mu(t),\qquad
 c_n=\int_0^\infty t^n\,d\nu(t).
\]
They have every positive exponential moment, unbounded support, and no
atom at zero. Define, for $s>0$,
\begin{equation}\label{gold:stieltjes}
 F(s)=\int\frac{d\mu(t)}{1+st},\qquad
 G(s)=\int\frac{d\nu(t)}{1+st},\qquad U(s)=sF(s).
\end{equation}
Then
\begin{equation}\label{gold:transform}
 F(s)=1-sF(s)^2F(sF(s)),\qquad
 G(s)=\frac1{1+U(s)},\qquad
 U^{-1}(s)=s(1+U(s)).
\end{equation}
Here $U$ is an increasing homeomorphism of $(0,\infty)$ onto itself.
With
\[
 \beta=\frac{\sqrt5-1}{2},\qquad
 \alpha=1-\beta=\beta^2,\qquad \varphi=\beta^{-1},
\]
the endpoint equivalents are
\begin{align}
 F(s)&\sim s^{-\alpha},&G(s)&\sim s^{-\beta} \quad(s\to\infty),
 \label{gold:leading}\\
 \mu([0,t])&\sim\frac{\sin(\pi\alpha)}{\pi\alpha}t^\alpha,&
 \nu([0,t])&\sim\frac{\sin(\pi\beta)}{\pi\beta}t^\beta
 \quad(t\downarrow0).\label{gold:cdf}
\end{align}
The coefficients are log-convex, and
\begin{equation}\label{gold:ratios}
 r_n:=\frac{a_n}{a_{n-1}}\sim\frac n{\log n}.
\end{equation}
For every integer $N\ge0$, the finite expansion of $F(s)$ has the exact
signed remainder
\begin{equation}\label{gold:finite}
 F(s)=\sum_{j=0}^{N}(-1)^ja_js^j
 +(-1)^{N+1}s^{N+1}\int\frac{t^{N+1}}{1+st}\,d\mu(t).
\end{equation}
\end{proposition}

\paragraph{Dependency and proof mechanism.}
This proposition packages established statements of the pinned research
report, rather than assigning them a new proof or a formal verification
status here. Its construction iterates a moment-preserving map: a positive
block operator gives $(1-zB(z))^{-1}$, and a positive full-Fock operator
turns a moment series $R$ into the solution $M=R(zM)$. Their composition
is $M=1+zM^2B(zM)$. Each iteration stabilizes one more coefficient.
Tightness and uniform integrability give the two limiting measures.
The proved coefficient growth makes their exponential generating functions
entire and hence their moment problems determinate. Weak limits of compact
approximants give \eqref{gold:transform} on the negative real axis.

The leading constant in \eqref{gold:leading} is also part of the prior
theorem. For orientation, put $g(x)=\log U(e^x)$ and
$d(x)=g(x)-\beta x$. The inverse identity implies
$d(x)=-\beta d(g(x))-\beta\varepsilon(x)$ with
$\varepsilon(x)=-\log(1-F(e^x))\to0$. Iteration to a compact interval,
followed by a contraction bound, proves $d(x)\to0$.
The Stieltjes Tauberian argument then gives \eqref{gold:cdf}.
Log-convexity and the coefficient logarithmic growth give
\eqref{gold:ratios}. Formula \eqref{gold:finite} follows directly from
the finite geometric identity and will make the new oscillation certificate
independent of an infinite numerical-series approximation.

The new proofs separate into three directions. Complex inversion improves
the integrated endpoint laws to analytic densities and derivative
asymptotics. Real inverse iteration identifies the missing oscillatory
correction. A composition estimate sharpens the coefficient comparison
without requiring the finer Bell-number conjecture.



% Included source: sections/golden_regularity.tex
\section{A088714: positive analytic densities on the full half-line}
\label{reg:section}

The moment representation in the repository leaves a concrete question:
can either of its two canonical measures have a singular component,
a positive atom, or a gap in its support?  We resolve all three issues.
The proof combines a global inversion theorem with a descent argument
specific to the coupled transforms.  A separate argument then determines
the densities at the origin; it does not differentiate the previously
known distribution-function equivalents.

Throughout this section, $\mu$ and $\nu$ are the canonical probability
measures on $[0,\infty)$ with moment series $A$ and
$C=(1-zA)^{-1}$, respectively.  We use their inherited finite moments,
absence of atoms at zero, and the negative-axis identity
\begin{equation}\label{reg:eq:inheritedU}
 s=U(s)\bigl(1+U(U(s))\bigr),\qquad
 U(s)=s\int_{[0,\infty)}\frac{d\mu(t)}{1+st}.
\end{equation}
For the endpoint calculation we also use the inherited law
\begin{equation}\label{reg:eq:inheritedCDF}
 \nu([0,t])\sim\frac{\sin(\pi\beta)}{\pi\beta}t^\beta,
 \qquad t\downarrow0,
 \qquad
 \beta=\frac{\sqrt5-1}{2},\quad
 \alpha=1-\beta=\beta^2.
\end{equation}
No assertion about the existence of a density is included in these
inputs.

\begin{theorem}[Regularity, full support, and density at zero]
\label{reg:thm:main}
Both $\mu$ and $\nu$ are absolutely continuous with respect to Lebesgue
measure.  Their densities $p$ and $q$ have strictly positive,
real-analytic versions on $(0,\infty)$, and
\begin{equation}\label{reg:eq:fullsupport}
 \operatorname{supp}\mu=\operatorname{supp}\nu=[0,\infty).
\end{equation}
In particular, neither law has atoms, a singular continuous component,
or a gap in its support.  Moreover,
\begin{equation}\label{reg:eq:densitieszero}
 p(x)\sim\frac{\sin(\pi\alpha)}{\pi}x^{-\beta},
 \qquad
 q(x)\sim\frac{\sin(\pi\beta)}{\pi}x^{-\alpha},
 \qquad x\downarrow0.
\end{equation}
The constants in the two density equivalents are equal because
$\alpha+\beta=1$.
\end{theorem}

\subsection{The analytic inverse selected by the moment laws}
\label{reg:subsec:transforms}

Write
\[
 \mathcal G_\mu(z)=\int\frac{d\mu(t)}{z-t},\qquad
 \mathcal G_\nu(z)=\int\frac{d\nu(t)}{z-t},\qquad
 K(z)=\frac1{\mathcal G_\mu(z)}.
\]
These functions are holomorphic on
$\mathbb C\setminus[0,\infty)$.  Each Cauchy transform maps the upper
half-plane into the lower half-plane, so $K$ maps the upper half-plane
into itself.  A real argument $x<0$ gives $K(x)<0$.

The renewal identity, initially justified on the negative axis, extends
by the identity theorem to give
\begin{equation}\label{reg:eq:renewalCauchy}
 \mathcal G_\nu(z)
   =\frac1{z\bigl(1-\mathcal G_\mu(z)\bigr)}.
\end{equation}
Define
\begin{equation}\label{reg:eq:H}
 H(w)=w^2\mathcal G_\nu(w)
      =\frac{w}{1-\mathcal G_\mu(w)}.
\end{equation}
Since $\nu$ is a probability measure with finite second moment,
polynomial division yields the integral representation
\begin{equation}\label{reg:eq:Hintegral}
 H(w)=w+m_1(\nu)+\int\frac{t^2}{w-t}\,d\nu(t).
\end{equation}
In the present normalization $m_1(\nu)=1$.

\begin{lemma}[The exact inverse identity]\label{reg:lem:inverse}
For every $z$ in the upper half-plane,
\begin{equation}\label{reg:eq:inverse}
 H(K(z))=z.
\end{equation}
\end{lemma}
\begin{proof}
For $s>0$ one has
$\mathcal G_\mu(-1/s)=-U(s)$ and $K(-1/s)=-1/U(s)$.
Consequently, \eqref{reg:eq:H} and \eqref{reg:eq:inheritedU} give
\[
 H(K(-1/s))
 =-\frac1{U(s)\bigl(1+U(U(s))\bigr)}
 =-\frac1s.
\]
The composition $H\circ K$ is holomorphic on
$\mathbb C\setminus[0,\infty)$: $K$ maps each half-plane into itself
and the negative real axis into itself.  The identity theorem now proves
\eqref{reg:eq:inverse}.
\end{proof}

Thus the required inverse relation follows from identities of actual
transforms.  It does not depend on evaluating the divergent series $A$
at a positive argument.  In free-probability terminology,
\eqref{reg:eq:Hintegral} identifies $\mu$ as the rate-one free compound
Poisson law with jump distribution $\nu$.  The argument below uses the
explicit analytic representation.

\subsection{The global inversion input and its hypotheses}
\label{reg:subsec:global}

We use the following special form of the global inversion theorem of
Huang and Wang \cite[Propositions 3.1--3.2]{HuangWang2022}.

\begin{theorem}[Global inversion; imported result]
\label{reg:thm:imported}
Let $\tau$ be a nonzero finite positive measure on $\mathbb R$, let
$b\in\mathbb R$, and put
\[
 H_\tau(w)=w+b+\int\frac{d\tau(t)}{w-t},\qquad
 I_\tau(u,v)=\int\frac{d\tau(t)}{(u-t)^2+v^2}.
\]
Define
\[
 v_\tau(u)=\inf\{v>0:I_\tau(u,v)<1\},\qquad
 \Omega_\tau=\{u+iv:v>v_\tau(u)\}.
\]
Then $v_\tau$ is finite and continuous, and
$v_\tau(u)=0$ exactly when
$\int (u-t)^{-2}\,d\tau(t)\le1$.
The map $H_\tau:\Omega_\tau\to\mathbb C^+$ is a holomorphic bijection.
It and its inverse extend continuously as mutually inverse maps between
their closures.  The boundary is
$\partial\Omega_\tau=\{u+iv_\tau(u):u\in\mathbb R\}$, and
\[
 h_\tau(u)=H_\tau(u+iv_\tau(u))
\]
is a strictly increasing homeomorphism of $\mathbb R$ onto itself.
For $v_\tau(u)>0$ one has
$I_\tau(u,v_\tau(u))=1$.
\end{theorem}

Here and below a boundary value of $H_\tau$ at a real point means its
continuous extension from $\Omega_\tau$.  The theorem is a direct
specialization of the cited Nevanlinna integral form: take
$d\sigma(t)=(1+t^2)^{-1}d\tau(t)$ and replace the real constant by
$b-\int t\,d\sigma(t)$.  All these integrals are finite because $\tau$
is finite.

For \eqref{reg:eq:Hintegral}, set
\begin{equation}\label{reg:eq:tau}
 d\tau(t)=t^2\,d\nu(t),\qquad b=m_1(\nu).
\end{equation}
The second moment of $\nu$ makes $\tau$ finite; its first moment is
positive, so $\tau\ne0$.  Thus every hypothesis of
Theorem~\ref{reg:thm:imported} holds.  Suppress the subscript $\tau$ and
write $v$, $\Omega$, and $h$ for the resulting objects.

For $z\in\mathbb C^+$, \eqref{reg:eq:inverse} implies
$\operatorname{Im}H(K(z))>0$.  The integral formula gives
\begin{equation}\label{reg:eq:imagH}
 \operatorname{Im}H(u+iy)
 =y\left(1-\int\frac{t^2}{(u-t)^2+y^2}\,d\nu(t)\right).
\end{equation}
It follows that $K(z)\in\Omega$.  Uniqueness of the inverse in
Theorem~\ref{reg:thm:imported} therefore identifies $K$ with its global
inverse.  In particular, $K$ has a continuous, injective extension to
$\mathbb C^+\cup\mathbb R$, and
\begin{equation}\label{reg:eq:boundaryinverse}
 K(h(u))=u+iv(u),\qquad \operatorname{Re}K(x)=h^{-1}(x).
\end{equation}

The absence of an atom of $\nu$ at zero gives
\[
 \int\frac{t^2}{(0-t)^2}\,d\nu(t)=1.
\]
Hence $v(0)=0$, while \eqref{reg:eq:Hintegral} gives the boundary value
$H(0)=m_1(\nu)-\int t\,d\nu(t)=0$.  The evaluation is absolutely
convergent, and it also follows by dominated convergence along the
positive imaginary axis.  Therefore
\begin{equation}\label{reg:eq:Ksign}
 h(0)=0,\qquad K(0)=0,\qquad
 K(x)\ne0\ (x\ne0),\qquad
 \operatorname{Re}K(x)>0\ (x>0).
\end{equation}

\subsection{Absolute continuity and the descent that excludes zeros}
\label{reg:subsec:descent}

On every compact interval in $(0,\infty)$, the reciprocal
$\mathcal G_\mu=1/K$ has finite, continuous boundary values.  This already
proves local absolute continuity.  Indeed, the Poisson integral of $\mu$
is $-\pi^{-1}\operatorname{Im}\mathcal G_\mu(x+iy)$.  On such an interval
it converges uniformly to a continuous function, whereas integration
against a compactly supported continuous test function converges to
integration against $\mu$.  The two limits identify the density as
\begin{equation}\label{reg:eq:densityK}
 p(x)=-\frac1\pi\operatorname{Im}\mathcal G_\mu(x)
      =\frac{\operatorname{Im}K(x)}{\pi|K(x)|^2},\qquad x>0.
\end{equation}
The inherited absence of an atom at zero, together with support in
$[0,\infty)$, proves that $\mu$ is globally absolutely continuous.
At this point $p$ is nonnegative and continuous; positivity still needs
proof.

\begin{lemma}[Descent of a hypothetical density zero]
\label{reg:lem:descent}
Let $Z=\{x>0:p(x)=0\}$.  If $x\in Z$, then
\begin{equation}\label{reg:eq:descent}
 w=K(x)\in Z,\qquad 0<w<x,\qquad K(w)>1.
\end{equation}
\end{lemma}
\begin{proof}
By \eqref{reg:eq:densityK}, $K(x)$ is real, and
\eqref{reg:eq:Ksign} makes $w=K(x)>0$.
Since $K(x)$ lies on $\partial\Omega$, one has $v(w)=0$.
Thus the vertical path $w+iy$, $y>0$, lies in $\Omega$, and the global
boundary inverse gives
\begin{equation}\label{reg:eq:Hwfinite}
 \lim_{y\downarrow0}H(w+iy)=H(w)=x.
\end{equation}
The value $\mathcal G_\mu(w)=1/K(w)$ is finite, since $w>0$ and
$K$ vanishes only at zero.  If
$1-\mathcal G_\mu(w)=0$, the quotient in \eqref{reg:eq:H} would have
unbounded modulus along this vertical path, contradicting
\eqref{reg:eq:Hwfinite}.  We may therefore pass to the boundary in that
quotient and obtain
\begin{equation}\label{reg:eq:zerotransfer}
 \mathcal G_\mu(w)=1-\frac wx.
\end{equation}
This is real, so \eqref{reg:eq:densityK} shows that $w\in Z$.
Applying \eqref{reg:eq:Ksign} at $w$ now gives
$\mathcal G_\mu(w)=1/K(w)>0$.  Equation
\eqref{reg:eq:zerotransfer} consequently gives $w<x$.
It also gives $\mathcal G_\mu(w)<1$, because $w>0$ and $x>0$.
Hence $K(w)>1$.
\end{proof}

\begin{proposition}[Every positive point has positive density]
\label{reg:prop:positive}
The set $Z$ is empty.
\end{proposition}
\begin{proof}
Suppose $x_0\in Z$ and define $x_{j+1}=K(x_j)$ for $j\ge0$.
Lemma~\ref{reg:lem:descent} makes this an infinite strictly decreasing
sequence of positive elements of $Z$.  Equation
\eqref{reg:eq:zerotransfer} gives
\[
 \frac1{x_{j+2}}=1-\frac{x_{j+1}}{x_j},\qquad
 x_j-x_{j+1}=\frac{x_j}{x_{j+2}}>1.
\]
It follows that $x_j<x_0-j$ for every $j\ge1$, contradicting positivity
once $j>x_0$.
\end{proof}

The feedback identity forces a decrease greater than one at every
hypothetical zero.  Merely knowing that the measures have mass
arbitrarily close to zero and unbounded support would not exclude
intermediate gaps.

\subsection{Analyticity and the second density}
\label{reg:subsec:analyticity}

For $x>0$, write $K(x)=u+iv$; the preceding proposition gives $v>0$.
This point lies in the open upper half-plane and on $\partial\Omega$,
so
\[
 \int\frac{t^2}{(u-t)^2+v^2}\,d\nu(t)=1.
\]
Differentiating \eqref{reg:eq:Hintegral} at this interior point and
using this equality gives
\begin{equation}\label{reg:eq:noncritical}
 \operatorname{Re}H'(u+iv)
 =2v^2\int\frac{t^2}{\bigl((u-t)^2+v^2\bigr)^2}\,d\nu(t)>0.
\end{equation}
Thus $H'(K(x))\ne0$.  The holomorphic inverse function theorem extends
$K$ across a neighborhood of $x$.  Its reciprocal extends there as
well, and \eqref{reg:eq:densityK} proves that $p$ is real analytic.

In \eqref{reg:eq:renewalCauchy}, the boundary denominator cannot vanish
at any $x>0$, because
$\operatorname{Im}\mathcal G_\mu(x)=-\pi p(x)<0$.
Consequently $\mathcal G_\nu$ has a finite holomorphic continuation
through every positive point.  The same Poisson-inversion argument
proves local absolute continuity of $\nu$ and the formula
\begin{equation}\label{reg:eq:densitybridge}
 q(x)=\frac{p(x)}{x\,|1-\mathcal G_\mu(x)|^2}>0,
 \qquad x>0.
\end{equation}
It is real analytic.  The absence of an atom at zero completes global
absolute continuity.  Strict positivity on every positive interval
proves \eqref{reg:eq:fullsupport} for both laws.

For completeness, the boundary parametrization gives exact simultaneous
density formulas.  For $u>0$, the number $v(u)$ is positive and uniquely
determined by
\begin{equation}\label{reg:eq:boundaryequation}
 \int\frac{t^2}{(u-t)^2+v(u)^2}\,d\nu(t)=1,
 \qquad h(u)=H(u+iv(u)).
\end{equation}
Equations \eqref{reg:eq:densityK} and \eqref{reg:eq:densitybridge} become
\begin{equation}\label{reg:eq:parametricdensities}
 \begin{split}
 p(h(u))&=\frac{v(u)}{\pi\bigl(u^2+v(u)^2\bigr)},\\
 q(h(u))&=\frac{v(u)}{
       \pi h(u)\bigl((u-1)^2+v(u)^2\bigr)}.
 \end{split}
\end{equation}
These identities are exact.  They characterize the density using the
same probability law $\nu$, so they are not by themselves a closed-form
numerical solution.

\subsection{A complex Abelian lemma at the origin}
\label{reg:subsec:abelian}

The transition from a distribution-function asymptotic to a density
asymptotic requires more than differentiation.  We first obtain a
uniform complex-transform limit away from the real support and then
invert $H$ at an interior point of the upper half-plane.

\begin{lemma}[Scaled Cauchy transforms]\label{reg:lem:abelian}
Let $\sigma$ be a probability measure on $[0,\infty)$, and suppose
\[
 N(t)=\sigma([0,t])\sim c\,t^\gamma\quad(t\downarrow0),
 \qquad c>0,\quad 0<\gamma<1.
\]
Then, locally uniformly for
$\zeta\in\mathbb C\setminus[0,\infty)$,
\begin{equation}\label{reg:eq:abelian}
 r^{1-\gamma}\mathcal G_\sigma(r\zeta)
 \longrightarrow
 -\frac{c\pi\gamma}{\sin(\pi\gamma)}(-\zeta)^{\gamma-1}
 \qquad(r\downarrow0).
\end{equation}
The power uses the branch of $\log(-\zeta)$ that is real for
$\zeta<0$.
\end{lemma}
\begin{proof}
The hypothesis implies $N(0)=0$ and a global bound
$N(t)\le C_0t^\gamma$ for a finite constant $C_0$.
Integration by parts, followed by $t=ry$, gives
\[
 \mathcal G_\sigma(z)
 =-\int_0^\infty\frac{N(t)}{(z-t)^2}\,dt,
 \qquad
 r^{1-\gamma}\mathcal G_\sigma(r\zeta)
 =-\int_0^\infty
       \frac{N(ry)/r^\gamma}{(\zeta-y)^2}\,dy.
\]
On each compact subset of the slit plane, the absolute value of the
integrand is bounded by a constant times
$y^\gamma/(1+y)^2$, which is integrable.  Dominated convergence is
therefore locally uniform and gives the limit
$-c\int_0^\infty y^\gamma(\zeta-y)^{-2}\,dy$.
For $\zeta=-a<0$, Euler's beta integral evaluates it as
\[
 -c\,a^{\gamma-1}
    \Gamma(\gamma+1)\Gamma(1-\gamma)
 =-\frac{c\pi\gamma}{\sin(\pi\gamma)}a^{\gamma-1}.
\]
Both expressions are holomorphic on the slit plane, so the identity
theorem evaluates the limit everywhere in that plane.
\end{proof}

Apply this lemma to \eqref{reg:eq:inheritedCDF}.  Since
$\alpha=1-\beta$, it gives, locally uniformly for $\zeta\in\mathbb C^+$,
\begin{equation}\label{reg:eq:nscaled}
 r^\alpha\mathcal G_\nu(r\zeta)
 \longrightarrow-(-\zeta)^{-\alpha}
 =e^{-i\pi\beta}\zeta^{-\alpha}.
\end{equation}
The last expression uses $0<\arg\zeta<\pi$.

\subsection{Inversion at a nonreal point gives the density constants}
\label{reg:subsec:endpoint}

\begin{proposition}[Boundary Cauchy transforms at zero]
\label{reg:prop:boundaryzero}
The continuous boundary values on $(0,\infty)$ satisfy
\begin{equation}\label{reg:eq:boundaryzero}
 K(x)\sim e^{i\pi\alpha}x^\beta,\qquad
 \mathcal G_\mu(x)\sim e^{-i\pi\alpha}x^{-\beta},\qquad
 \mathcal G_\nu(x)\sim e^{-i\pi\beta}x^{-\alpha}
 \qquad(x\downarrow0).
\end{equation}
Here a complex equivalent means that the quotient tends to $1$.
\end{proposition}
\begin{proof}
Use $H(w)=w^2\mathcal G_\nu(w)$, take $r=x^\beta$ in
\eqref{reg:eq:nscaled}, and note that
$\beta(2-\alpha)=\beta(1+\beta)=1$.  Thus
\begin{equation}\label{reg:eq:Hscaled}
 \frac{H(x^\beta\zeta)}x
 \longrightarrow e^{-i\pi\beta}\zeta^{1+\beta}
\end{equation}
locally uniformly in $\mathbb C^+$.
The equation
\[
 e^{-i\pi\beta}\zeta^{1+\beta}=1
\]
has a simple root at
$\zeta_0=e^{i\pi\alpha}$, because
$\alpha(1+\beta)=\beta$.

Choose an arbitrarily small closed disk about $\zeta_0$ contained
strictly in the upper half-plane and containing no other zero of
$e^{-i\pi\beta}\zeta^{1+\beta}-1$.
Uniform convergence on its boundary and Rouch\'e's theorem show that,
for all sufficiently small $x>0$, the equation
$H(x^\beta\zeta)=x$ has a root in this disk.
Let $w=x^\beta\zeta$ be that root.
It lies in $\mathbb C^+$ and has $\operatorname{Im}H(w)=0$.
By \eqref{reg:eq:imagH} and strict monotonicity of the integral there
as a function of the positive imaginary part, $w$ lies on
$\partial\Omega$.
The boundary bijection in Theorem~\ref{reg:thm:imported} now forces
$w=K(x)$.
Since the disk can be arbitrarily small,
$K(x)/x^\beta\to e^{i\pi\alpha}$.

Taking reciprocals proves the second equivalent in
\eqref{reg:eq:boundaryzero}.  Finally,
\eqref{reg:eq:renewalCauchy} and
$|\mathcal G_\mu(x)|\to\infty$ give
\[
 \mathcal G_\nu(x)
 \sim-\frac1{x\mathcal G_\mu(x)}
 \sim-e^{i\pi\alpha}x^{\beta-1}
 =e^{-i\pi\beta}x^{-\alpha}.
\]
\end{proof}

Taking negative imaginary parts of
\eqref{reg:eq:boundaryzero} proves \eqref{reg:eq:densitieszero} and
finishes Theorem~\ref{reg:thm:main}.  The proof uses a complex disk a
fixed positive distance from the real axis after scaling.  It does not
assert uniformity of an Abelian theorem up to the support, where such
uniformity would require an additional argument.

\begin{corollary}[Every fixed derivative at the origin]
\label{reg:cor:derivatives}
For each fixed integer $j\ge0$,
\begin{equation}\label{reg:eq:derivatives}
 \begin{split}
 p^{(j)}(x)&\sim
   \frac{\sin(\pi\alpha)}\pi
   (-1)^j(\beta)_j\,x^{-\beta-j},\\
 q^{(j)}(x)&\sim
   \frac{\sin(\pi\beta)}\pi
   (-1)^j(\alpha)_j\,x^{-\alpha-j},
 \qquad x\downarrow0,
 \end{split}
\end{equation}
where $(a)_0=1$ and $(a)_j=a(a+1)\cdots(a+j-1)$ for $j\ge1$.
In particular, each density is strictly decreasing and strictly convex
on a sufficiently short interval adjacent to zero.
\end{corollary}
\begin{proof}
The same Rouch\'e argument may be made uniform when the right-hand side
$1$ in the scaled equation is replaced by a complex variable $z$ in a
small disk about $1$.
Choose the disk small enough that
$e^{i\pi\alpha}z^\beta$ stays in a fixed disk compactly contained in
$\mathbb C^+$.  The limiting inverse is holomorphic with nonzero
derivative.  Local uniform convergence of the scaled $H$ also gives
local uniform convergence of its derivatives by Cauchy's formula, so
the nearby roots remain noncritical and define holomorphic inverses by
the implicit function theorem.  Uniqueness identifies these inverses with
$x^{-\beta}K(xz)$ for $\operatorname{Im}z>0$, and then with their
continuations through the positive axis.  They converge locally
uniformly to $e^{i\pi\alpha}z^\beta$.
Consequently
\[
 x^\beta\mathcal G_\mu(xz)
   \longrightarrow e^{-i\pi\alpha}z^{-\beta},\qquad
 x^\alpha\mathcal G_\nu(xz)
   \longrightarrow e^{-i\pi\beta}z^{-\alpha}
\]
locally uniformly in that disk.  Cauchy's differentiation formula
permits any fixed number of derivatives at $z=1$.
Taking negative imaginary parts gives \eqref{reg:eq:derivatives}.
The cases $j=1,2$ give the final assertion.
\end{proof}

\subsection{What this resolves and what remains}
\label{reg:subsec:scope}

Theorem~\ref{reg:thm:main} resolves the repository's question about
regularity and support of both canonical laws.  The additional boundary
argument provides the density equivalents that were expressly left
unjustified by the earlier cumulative-mass estimates.  The general
global inversion theorem is an imported tool; the exclusion of every
positive density zero comes from the particular feedback identity
$H(w)=w/(1-\mathcal G_\mu(w))$ and
Lemma~\ref{reg:lem:descent}.

These results leave substantive questions.  They do not give a
multiplicative equivalent for the far tails or a complete large-$n$
asymptotic for the moments.  Nor does local decrease near zero prove
global monotonicity of either density.  Useful next problems are to
determine whether $p$ or $q$ is globally decreasing, to obtain
quantitative large-$x$ boundary estimates from \eqref{reg:eq:H}, and to
turn \eqref{reg:eq:boundaryequation} into certified density enclosures
with a proved convergence rate.  The derivative equivalents are
statements for each fixed order $j$; they do not establish complete
monotonicity on one common interval for all orders.



% Included source: sections/golden_dynamics.tex
\section{A complete oscillatory expansion for the golden-ratio laws}
\label{dyn:sec:main}

The repository asks whether the leading golden-ratio asymptotics for
A088714 admit an explicit second-order expansion, and whether the
correction oscillates along inverse orbits. This section answers both
questions. The correction is of order $1/\log s$, has a nonzero
real-analytic profile periodic in $\log\log s$, and admits a convergent
expansion in inverse powers of $\log s$, with an error smaller than all
those powers. The profile changes sign after one
golden-ratio dilation of the logarithmic variable. An exact interval
certificate proves that this oscillation is present in the particular
OEIS solution, rather than merely permitted by the functional equation.

\subsection{Inherited input and the main statement}
\label{dyn:sub:input}

We use the following results from Part~II of the repository manuscript
\cite{goldenprior}: the theorems entitled \emph{Moment and analytic
structure} and \emph{Golden-ratio endpoint laws}, together with its
Stieltjes remainder formula. Its source labels are
\texttt{gml:thm:mainmoment}, \texttt{gml:thm:mainendpoint}, and
\texttt{gml:eq:remainder}. In the notation of that manuscript,
\begin{equation}\label{dyn:eq:input}
 F(s)=\int_{[0,\infty)}\frac{d\mu(t)}{1+st},\qquad
 U(s)=sF(s),\qquad G(s)=\frac{1}{1+U(s)}.
\end{equation}
Here $\mu$ is the probability measure with moments A088714, $G$ is the
Stieltjes transform for A088713, $U$ is an increasing homeomorphism of
$(0,\infty)$ onto itself, and
\begin{equation}\label{dyn:eq:inverse}
 U^{-1}(t)=t(1+U(t)),\qquad
 U(s)\sim s^\beta,
 \qquad \beta=\frac{\sqrt5-1}{2}.
\end{equation}
We put
\begin{equation}\label{dyn:eq:constants}
 \varphi=\frac{1+\sqrt5}{2}=\frac1\beta,\qquad
 \alpha=1-\beta=\beta^2,\qquad \lambda=\log\varphi,
 \qquad \mathcal D=\lambda^{-1}\frac{d}{dt}.
\end{equation}
The letter $P$ below denotes an analytic function, not a polynomial. Its
definition is constructive and canonical; it is not a freely chosen
periodic multiplier.

\begin{theorem}[Complete golden-ratio expansion]\label{dyn:thm:allorders}
There is a nonzero real-valued real-analytic function $P$ on $\mathbb R$
such that
\begin{equation}\label{dyn:eq:anti}
 P(t+1)=-P(t),\qquad P(t+2)=P(t).
\end{equation}
For every fixed integer $N\ge1$, as $s\to\infty$,
\begin{align}
 \log\bigl(s^\alpha F(s)\bigr)
 &=\sum_{n=1}^{N}
   \frac{Q_n(\log_\varphi\log s)}{(\log s)^{2n-1}}
   +O\bigl((\log s)^{-2N-1}\bigr),\label{dyn:eq:Ffull}\\
 \log\bigl(s^\beta G(s)\bigr)
 &=-\sum_{n=1}^{N}
   \frac{Q_n(\log_\varphi\log s)}{(\log s)^{2n-1}}
   +O\bigl((\log s)^{-2N-1}\bigr),\label{dyn:eq:Gfull}
\end{align}
where an empty product is $1$ and
\begin{equation}\label{dyn:eq:Qformula}
 Q_n(t)=\frac{5^{(1-n)/2}}{n!}
 \left[\prod_{j=n}^{2n-2}(\mathcal D-j)\right]P(t)^n.
\end{equation}
The remainders hold uniformly over the oscillation phase. Moreover, for
all sufficiently large $x$, the series
\begin{equation}\label{dyn:eq:Sconvergent}
 S(x)=\sum_{n=1}^{\infty}\frac{Q_n(\log_\varphi x)}{x^{2n-1}}
\end{equation}
converges absolutely, uniformly over the phase, and
\begin{align}
 \log(s^\alpha F(s))&=S(\log s)+O(s^{-\beta}),
       \label{dyn:eq:Fsector}\\
 \log(s^\beta G(s))&=-S(\log s)+O(s^{-1}).
       \label{dyn:eq:Gsector}
\end{align}
In particular a first correction beyond all inverse-logarithmic orders
can be isolated explicitly:
\begin{equation}\label{dyn:eq:beyond}
 s^\alpha F(s)=\exp(S(\log s))-s^{-\beta}+O(s^{-1}).
\end{equation}
\end{theorem}

For example,
\begin{equation}\label{dyn:eq:Qfirst}
 Q_1=P,\qquad
 Q_2=\frac{P\mathcal DP-P^2}{\sqrt5},\qquad
 Q_3=\frac1{30}(\mathcal D-3)(\mathcal D-4)P^3.
\end{equation}
In general $Q_n(t+1)=(-1)^nQ_n(t)$: the odd-indexed coefficient
functions are antiperiodic and the even-indexed ones are $1$-periodic.
Thus, writing $x=\log s$ and $p=P(\log_\varphi x)$,
\begin{align}
 s^\alpha F(s)
 &=1+\frac{p}{x}+\frac{p^2}{2x^2}+O(x^{-3}),\label{dyn:eq:Ffirst}\\
 s^\beta G(s)
 &=1-\frac{p}{x}+\frac{p^2}{2x^2}+O(x^{-3}).\label{dyn:eq:Gfirst}
\end{align}
There is no constant coefficient of $1/\log s$: the leading correction
is the nonconstant sign-changing function $P$ evaluated at
$\log_\varphi\log s$.

\subsection{The exact perturbed Fibonacci recurrence}
\label{dyn:sub:recurrence}

Set
\begin{equation}\label{dyn:eq:gdef}
 g(x)=\log U(e^x),\qquad h=g^{-1},\qquad
 d(x)=g(x)-\beta x,\qquad r(y)=\log(1+e^{-y}).
\end{equation}
The inherited equivalent in \eqref{dyn:eq:inverse} gives $d(x)\to0$.
The inverse identity becomes
\begin{equation}\label{dyn:eq:h}
 h(x)=x+\log(1+e^{g(x)})=x+g(x)+r(g(x)).
\end{equation}
Fix $x$ sufficiently large and write
\begin{equation}\label{dyn:eq:orbit}
 x_{-1}=g(x),\quad x_0=x,\quad x_n=h^n(x)\quad(n\ge0).
\end{equation}
Then, exactly,
\begin{equation}\label{dyn:eq:fibonacci}
 x_{n+1}=x_n+x_{n-1}+r(x_{n-1})\qquad(n\ge0).
\end{equation}
In particular the forcing is positive and becomes smaller than every
geometrically decreasing sequence along a positive orbit. If $f_n$ are
the Fibonacci numbers, $f_0=0$, $f_1=1$, induction gives
\begin{equation}\label{dyn:eq:fiblower}
 x_n\ge f_{n+1}x+f_ng(x)>0\qquad(n\ge0).
\end{equation}
On every compact interval of initial points with $g(x)>0$, the right
side is bounded below by $c\varphi^n$ with some $c>0$.

Define $d_n=x_{n-1}-\beta x_n=d(x_n)$. The characteristic roots of the
homogeneous recurrence are $\varphi$ and $-\beta$, and direct
substitution gives the stable-coordinate recurrence
\begin{equation}\label{dyn:eq:stable}
 d_{n+1}=-\beta d_n-\beta r(x_{n-1}).
\end{equation}
The negative sign is the exact source of the alternating correction.

\begin{proposition}[Canonical growing and alternating coordinates]
\label{dyn:prop:coordinates}
On a sufficiently large ray, the following limits and series exist:
\begin{align}
 \Psi(x)
 &=\lim_{n\to\infty}\beta^n h^n(x)
  =x+\frac{d(x)}{\sqrt5}
    +\frac1{\sqrt5}\sum_{j=0}^{\infty}
      \beta^j r(x_{j-1}),\label{dyn:eq:Psi}\\
 D(x)
 &=\lim_{n\to\infty}(-\beta)^{-n}d(h^n(x))
  =d(x)+\sum_{j=0}^{\infty}(-\varphi)^j r(x_{j-1}).
  \label{dyn:eq:D}
\end{align}
Both series converge absolutely and locally uniformly. Moreover,
\begin{equation}\label{dyn:eq:conjugacies}
 \Psi(h(x))=\varphi\Psi(x),\qquad D(h(x))=-\beta D(x).
\end{equation}
The function $\Psi$ is positive, real analytic and strictly increasing, has
$\Psi'(x)>0$, and maps any sufficiently large closed ray
$[X,\infty)$ homeomorphically onto $[\Psi(X),\infty)$.
\end{proposition}

\begin{proof}
Iteration of \eqref{dyn:eq:stable} proves the formula for $D$.
Variation of constants in \eqref{dyn:eq:fibonacci} gives
\[
 x_n=f_{n+1}x+f_ng(x)
       +\sum_{j=0}^{n-1}f_{n-j}r(x_{j-1}).
\]
Insert Binet's formula
$f_n=(\varphi^n-(-\beta)^n)/\sqrt5$ and multiply by $\beta^n$.
The bound \eqref{dyn:eq:fiblower}, together with $0<r(y)\le e^{-y}$ for
$y>0$, justifies the limit and yields \eqref{dyn:eq:Psi}. Absolute and
locally uniform convergence of both series follows from the same bound.
Shifting the limits by one iterate proves
\eqref{dyn:eq:conjugacies}. Positivity of $\Psi$ follows also from its
equivalent formula with initial term
$(\varphi x+g(x))/\sqrt5$ and positive summands.

We prove analyticity rather than assuming regularity of the limiting
coordinate. The Stieltjes representation extends $U$ holomorphically
to a neighborhood of every positive real point. Since $U$ is positive
there, $g$ is holomorphic in a neighborhood of every real point. On the
real axis, differentiation under the integral, with
$w_t(x)=e^x/(1+te^x)$, gives
\begin{equation}\label{dyn:eq:gprime}
 0<g'(x)=
 \frac{\int w_t(x)(1+te^x)^{-1}\,d\mu(t)}
      {\int w_t(x)\,d\mu(t)}<1.
\end{equation}
The strict upper bound follows because $\mu$ is not supported at zero.
In particular $h'>0$.

Fix a real initial point far enough to the right that it and its
$g$-image exceed $2$. On a sufficiently small complex disk around that
point, both $\operatorname{Re}z$ and $\operatorname{Re}g(z)$ exceed $2$.
Set $z_{-1}=g(z)$, $z_0=z$, and define the complex orbit by the recurrence
\[
 z_{n+1}=z_n+z_{n-1}+r(z_{n-1}).
\]
Here $r(w)=\log(1+e^{-w})$ is holomorphic throughout
$\operatorname{Re}w>0$, by its convergent logarithm series. For
$y=\operatorname{Re}w\ge2$ it satisfies
\[
 |r(w)|\le-\log(1-e^{-y})<\frac y2.
\]
Thus, inductively, all real parts stay above $2$ and
\[
 \operatorname{Re}z_{n+1}\ge
 \operatorname{Re}z_n+\tfrac12\operatorname{Re}z_{n-1}.
\]
They grow at least like $c\kappa^n$, uniformly on the disk, where
$\kappa=(1+\sqrt3)/2>1$. Consequently the complex versions of both
invariant series converge locally uniformly and absolutely; even the
larger summands are bounded by $C\varphi^j\exp(-c\kappa^j)$.
They therefore define holomorphic functions agreeing with $\Psi,D$
on the real interval. The complex recurrence supplies the extension:
we do not require the principal expression for $g$ or $h$ to be defined
at every later complex iterate. The same Binet formula makes
$\beta^nz_n$ converge locally uniformly to $\Psi$, so its derivatives
also converge locally uniformly.

Strict positivity of $\Psi'$ requires an additional argument. Put
$p_n=dx_n/dx>0$ and
$c_n=(1+e^{-x_{n-1}})^{-1}$. Differentiating the recurrence gives
$p_{n+1}=p_n+c_np_{n-1}$. For $q_n=p_n+\beta p_{n-1}$,
\[
 q_{n+1}=\varphi q_n-(1-c_n)p_{n-1}
          \ge\varphi c_nq_n,
\]
because $p_{n-1}\le q_n/\beta$ and $\varphi\beta=1$. It follows that
\[
 (1+\beta^2)\Psi'(x)
 =\lim_{n\to\infty}\beta^nq_n
 \ge (1+\beta g'(x))\prod_{j=0}^{\infty}c_j>0.
\]
The infinite product is positive since
$\sum_j(1-c_j)\le\sum_j e^{-x_{j-1}}<\infty$.
Finally, the series for $\Psi$ and $d(x)\to0$ show that
$\Psi(x)=x+o(1)$. Strict increase and this equivalent establish the
claimed range and homeomorphism property.
\end{proof}

\subsection{Effective remainder bounds and the profile}
\label{dyn:sub:profile}

The series above also supply explicit error control. Suppose that
$a=g(x)>\log\varphi$ and $x>a$. For $j\ge1$ the recurrence implies
$x_{j-1}\ge x+(j-1)a$. Consequently
\begin{align}
 |D(x)-d(x)|
 &\le e^{-a}+\frac{\varphi e^{-x}}{1-\varphi e^{-a}},
 \label{dyn:eq:Derror}\\
 \left|\Psi(x)-x-\frac{D(x)}{\sqrt5}\right|
 &\le\frac{e^{-x}}{\sqrt5}
  \left(\frac{\varphi}{1-\varphi e^{-a}}
       +\frac{\beta}{1-\beta e^{-a}}\right).
 \label{dyn:eq:Psisharp}
\end{align}
The second bound uses the cancellation of the $j=0$ terms in
\eqref{dyn:eq:Psi} and \eqref{dyn:eq:D}. More explicitly, the expression
inside the absolute value is
\begin{equation}\label{dyn:eq:Eseries}
 E(x)=\frac1{\sqrt5}\sum_{j=1}^{\infty}
          \bigl(\beta^j-(-\varphi)^j\bigr)r(x_{j-1}).
\end{equation}
For truncation after the terms $j=0,\ldots,J-1$, $J\ge1$, the two
series have respective tail bounds
\begin{align}
 \left|\sum_{j=J}^{\infty}(-\varphi)^jr(x_{j-1})\right|
 &\le\frac{\varphi^J e^{-x_{J-1}}}
               {1-\varphi e^{-x_{J-2}}},\label{dyn:eq:Dtail}\\
 0\le\frac1{\sqrt5}\sum_{j=J}^{\infty}\beta^jr(x_{j-1})
 &\le\frac{\beta^J e^{-x_{J-1}}}
       {\sqrt5\bigl(1-\beta e^{-x_{J-2}}\bigr)}.
       \label{dyn:eq:Psitail}
\end{align}
The denominators are positive on the ray just specified. These bounds
follow by using $x_{J-2}$ as a lower bound for each subsequent orbit
increment. They require no estimate of an unspecified asymptotic
constant.

Define $P$ first on a sufficiently large ray in its argument by
\begin{equation}\label{dyn:eq:Pdefine}
 P(\log_\varphi\Psi(x))=\Psi(x)D(x).
\end{equation}
Proposition~\ref{dyn:prop:coordinates}, including $\Psi'>0$, makes this
well-defined and real analytic by the holomorphic inverse-function
theorem.
The two conjugacy identities immediately give
\[
 P(\log_\varphi\Psi(x)+1)=-P(\log_\varphi\Psi(x)).
\]
This identity extends $P$ uniquely to a real-analytic function on all of
$\mathbb R$ satisfying \eqref{dyn:eq:anti}. In particular $P$ and all
its derivatives are bounded. The normalization
$\Psi(x)/x\to1$ fixes the phase coordinate used here.

Since $\Psi(x)\sim x$, boundedness of $P$ gives $D(x)=O(x^{-1})$.
The inherited relation $g(x)=\beta x+o(1)$ and
\eqref{dyn:eq:Derror} then yield
\begin{equation}\label{dyn:eq:sharpcoords}
 \begin{aligned}
 d(x)&=\frac{P(\log_\varphi\Psi(x))}{\Psi(x)}+O(e^{-\beta x}),\\
 \Psi(x)&=x+\frac{D(x)}{\sqrt5}+O(e^{-x})=x+O(x^{-1}).
 \end{aligned}
\end{equation}
Thus an exact change of variable makes the leading profile accurate up
to an exponentially small error in $x=\log s$. The preceding displayed
bounds give explicit versions of both errors.

The companion transform has a slightly stronger error in this same
coordinate. Its defining identity gives
\[
 \log(e^{\beta x}G(e^x))=-d(x)-r(g(x)).
\]
The $j=0$ term in the series for $D-d$ cancels, so
\begin{equation}\label{dyn:eq:Gcoord}
 \left|\log(e^{\beta x}G(e^x))+D(x)\right|
 \le\frac{\varphi e^{-x}}{1-\varphi e^{-g(x)}}.
\end{equation}
The coordinate error for the logarithm of the normalized $F$ is
$O(s^{-\beta})$; for the logarithm of the normalized $G$ it is
$O(s^{-1})$.

\subsection{Proof of the expansion to every order}
\label{dyn:sub:reversion}

We now prove Theorem~\ref{dyn:thm:allorders}, except for nonvanishing,
which is certified in the next subsection. Let $y=\Psi(x)$,
$t=\log_\varphi x$, and $v=y/x-1$. Equations
\eqref{dyn:eq:Pdefine} and \eqref{dyn:eq:Eseries} give
\begin{equation}\label{dyn:eq:scaledimplicit}
 v=\eta H_t(v)+\frac{E(x)}x,\qquad
 \eta=\frac1{\sqrt5\,x^2},\qquad
 H_t(v)=\frac{P(t+\log_\varphi(1+v))}{1+v}.
\end{equation}
Here $v=O(x^{-2})$ and $E(x)=O(e^{-x})$. Real analyticity and periodicity
of $P$ give a holomorphic extension to some fixed strip around the real
phase axis: local extensions glue by uniqueness, and compactness of a
period provides a common positive width. The analytic implicit-function
theorem therefore produces a unique small solution
$v_0=v_0(\eta,t)$ of $v_0=\eta H_t(v_0)$, holomorphic in $\eta$ near
zero with a convergence radius uniform for real $t\bmod2$. Since
$|\eta\partial_vH_t(v)|\le1/2$ there, subtracting the two
equations shows
\begin{equation}\label{dyn:eq:vcompare}
 v-v_0=O(e^{-x}/x).
\end{equation}
For every fixed $N$, Taylor's theorem in $\eta$, uniformly on the phase
circle, gives the Lagrange expansion
\begin{equation}\label{dyn:eq:lagrange}
 v_0(\eta,t)=\sum_{n=1}^{N}\frac{\eta^n}{n!}
 \left.\frac{d^{n-1}}{dv^{n-1}}H_t(v)^n\right|_{v=0}
 +O(\eta^{N+1}).
\end{equation}
The coefficient identity can be established at the level of finite
Taylor jets: substitution into $v_0=\eta H_t(v_0)$ determines the
coefficients recursively and gives the displayed Lagrange coefficient.
The analyticity just proved additionally makes the infinite Taylor
series converge absolutely, uniformly for real phase and for $\eta$ in
a sufficiently small closed disk.
For an explicit majorant, choose a common circle $|v|=\rho$ in the
domain of $H_t$, and let $|H_t(v)|\le M_H$ there for every real phase.
Cauchy's coefficient bound gives
\[
 \bigl|[\eta^n]v_0(\eta,t)\bigr|
 =\frac1n\bigl|[v^{n-1}]H_t(v)^n\bigr|
 \le\frac{M_H^n}{n\rho^{n-1}}.
\]
In particular the increasing derivative order in
\eqref{dyn:eq:Qformula} introduces no hidden factorial divergence.

Writing $u=1+v$, each differentiation of
$u^{-n}P(t+\log_\varphi u)^n$ replaces its current power $u^{-j}$ by
$u^{-j-1}$ and applies $\mathcal D-j$ to the function of phase. Thus
\begin{equation}\label{dyn:eq:jet}
 \left.\frac{d^{n-1}}{dv^{n-1}}H_t(v)^n\right|_{v=0}
 =\left[\prod_{j=n}^{2n-2}(\mathcal D-j)\right]P(t)^n.
\end{equation}
Finally,
$d(x)=D(x)+O(e^{-\beta x})
      =\sqrt5\,xv_0+O(e^{-\beta x})$.
Multiply \eqref{dyn:eq:lagrange} by $\sqrt5x$ and use
\eqref{dyn:eq:jet}. This gives \eqref{dyn:eq:Ffull} with the remainder
$O(x\eta^{N+1})=O(x^{-2N-1})$ and the coefficients
\eqref{dyn:eq:Qformula}. Formula \eqref{dyn:eq:Gcoord} gives
\eqref{dyn:eq:Gfull}. Absolute convergence of the full Taylor series
proves \eqref{dyn:eq:Sconvergent}. Equations
\eqref{dyn:eq:vcompare} and \eqref{dyn:eq:Psisharp} give, more precisely,
$D(x)=\sqrt5xv_0+O(e^{-x})=S(x)+O(e^{-x})$. Combine this with
\eqref{dyn:eq:Derror} and \eqref{dyn:eq:Gcoord} to obtain
\eqref{dyn:eq:Fsector}--\eqref{dyn:eq:Gsector}. Finally the exact
identity
\[
 s^\alpha F(s)+s^{-\beta}=(s^\beta G(s))^{-1}
\]
and \eqref{dyn:eq:Gsector} prove \eqref{dyn:eq:beyond}, since $S(x)$
is bounded and tends to zero. This proves every asserted expansion,
without assuming differentiable asymptotics for $F$ or $G$ in advance.

\subsection{A certified nonzero profile value}
\label{dyn:sub:certificate}

The nonvanishing assertion is a finite computation followed by a
rigorous infinite-tail estimate. It can be reproduced from the source
code in this package; no externally supplied decimal value of $F$ is
needed.

Generate the exact integers $a_n$ by the inherited triangular recurrence
\begin{equation}\label{dyn:eq:momentrec}
 \begin{aligned}
 c_{0,k}&=1,\qquad c_{n,0}=0\quad(n\ge1),\qquad a_n=c_{n,1},\\
 c_{n,k}&=c_{n,k-1}+\sum_{i=0}^{n-1}c_{i,k}c_{n-1-i,k+1}.
 \end{aligned}
\end{equation}
With $s_0=1/100$, the exact Stieltjes remainder gives rational bounds
\begin{equation}\label{dyn:eq:seed}
 \sum_{j=0}^{31}a_j(-s_0)^j
 \le F(s_0)\le
 \sum_{j=0}^{30}a_j(-s_0)^j.
\end{equation}
Set $s_{-1}=s_0F(s_0)$ and
$s_{k+1}=s_k(1+s_{k-1})$. The inverse identity proves inductively that
$U(s_k)=s_{k-1}$. Hence the true, generally unknown seed in
\eqref{dyn:eq:seed} defines an orbit of the actual transform. Put
$x_*=\log s_{115}$, so $g(x_*)=\log s_{114}$.

The included script \texttt{code/verify\_dynamics.py} uses rational
arithmetic for \eqref{dyn:eq:momentrec} and \eqref{dyn:eq:seed}, then
100-digit Decimal intervals for the orbit and invariant series. Basic
arithmetic is rounded outwards. Each correctly rounded exponential,
logarithm and square root is enlarged to the adjacent representable
number on both sides, using the documented Decimal operations
\cite{python-decimal}. The full endpoints are recorded in
\texttt{data/dynamics\_certificate.json}. The following widened
decimal enclosures are among the script's exact assertions:
\begin{align}
 26.4807143957&<x_*<26.4807143959,\label{dyn:eq:certx}\\
 16.366&<g(x_*)<16.367,\label{dyn:eq:certg}\\
 0.00003580&<d(x_*)<0.00003582,\label{dyn:eq:certd}\\
 |D(x_*)-d(x_*)|&<0.000000079.\label{dyn:eq:certtail}
\end{align}
The last inequality is the explicit bound
\eqref{dyn:eq:Derror}. In particular $D(x_*)>0$, proving $P\not\equiv0$.
For a stronger numerical enclosure, the script takes the first five
terms of both invariant series and bounds the remaining terms by
\eqref{dyn:eq:Dtail}--\eqref{dyn:eq:Psitail}. It proves
\begin{equation}\label{dyn:eq:certP}
 0.00095032738531560
 <P(\log_\varphi\Psi(x_*))
 <0.00095032738531562.
\end{equation}
This interval encloses one value of $P$, not its global maximum.

Every decimal statement just displayed concerns the same true orbit
point admitted by the exact rational seed bounds. It does not claim an
enclosure for all points in the small interval
\eqref{dyn:eq:certx}. This distinction is essential when transporting a
functional identity from an interval-valued seed.

\begin{corollary}[Sharp rate and genuine alternation]
\label{dyn:cor:sharp}
Let $M=\max_{t\in[0,2]}|P(t)|$. Then $M>0.00095032738531560$ and
\begin{align}
 \limsup_{s\to\infty}(\log s)(s^\alpha F(s)-1)&=M,
 &\liminf_{s\to\infty}(\log s)(s^\alpha F(s)-1)&=-M,\label{dyn:eq:limsupF}\\
 \limsup_{s\to\infty}(\log s)(s^\beta G(s)-1)&=M,
 &\liminf_{s\to\infty}(\log s)(s^\beta G(s)-1)&=-M.\label{dyn:eq:limsupG}
\end{align}
Thus the $O(1/\log s)$ relative error is optimal: neither error is
$o(1/\log s)$, and neither is $O(s^{-\delta})$ for any $\delta>0$.
Along the inverse orbit starting at $s_{115}$,
\begin{equation}\label{dyn:eq:orbitlimit}
 \lim_{n\to\infty}(-1)^n(\log s_{115+n})
       \bigl(s_{115+n}^{\alpha}F(s_{115+n})-1\bigr)
 =P(\log_\varphi\Psi(x_*))>0.
\end{equation}
Both normalized transforms cross their leading power laws infinitely
often.
\end{corollary}

\begin{proof}
Equations \eqref{dyn:eq:Ffirst}--\eqref{dyn:eq:Gfirst} show that the
two scaled errors have limiting profiles $P$ and $-P$. Their arguments
$\log_\varphi\log s$ pass through every phase infinitely often.
Antiperiodicity makes the maximum and minimum opposite, giving the four
limits. On the specified orbit,
$\Psi(h^n(x_*))=\varphi^n\Psi(x_*)$ and
$D(h^n(x_*))=(-\beta)^nD(x_*)$ exactly. Since
$h^n(x_*)/\Psi(h^n(x_*))\to1$, the coordinate estimate
\eqref{dyn:eq:sharpcoords} proves \eqref{dyn:eq:orbitlimit}.
Eventual alternating signs on this orbit, followed by continuity, prove
the crossing assertions. The other claims follow from $M>0$.
\end{proof}

\subsection{Scope and further questions}
\label{dyn:sub:questions}

This resolves the repository's second-order golden-ratio question,
including its suggested orbit oscillation, and gives a convergent
inverse-logarithmic expansion with smaller power corrections. It is a
result about the real positive
Stieltjes-transform axis. It does not by itself prove the unresolved
fine Bell-number normalization of the coefficients, nor a second-order
distribution-function or density asymptotic for the measures.

Several concrete follow-up problems remain.
\begin{enumerate}
\item Determine the maximum $M$ with certified bounds over an entire
phase interval, and count the zeros of $P$ in a period. The certificate
above proves only one nonzero value and the existence of zeros forced
by antiperiodicity.
\item Locate the complex singularities of the analytic continuation of
$P$, and obtain explicit widths for its holomorphic strip. This would
make the convergence radius and constants in the inverse-logarithmic
expansion effective.
\item Develop certified Fourier bounds for $P$. Antiperiodicity permits
only odd frequencies in a period-two Fourier representation, and
analyticity gives exponential decay. Effective decay rates and bounds
on the individual coefficients remain to be determined.
\item Recover sharp second-order behavior of the representing measures
near zero. A suitable Tauberian theorem must preserve the slowly
varying oscillation and control its remainder; differentiation of the
leading distribution-function equivalent alone is insufficient.
\item Use the explicit invariant series to design efficient evaluation
over all phases. The forward inverse recurrence is especially accurate
once $g(x)>0$; obtaining uniformly efficient certified seeds across a
fundamental phase interval is a separate numerical problem.
\item Investigate analogous profiles for the family of nonlinear
inverse relations $U^{-1}(t)=t(1+U(t))^q$. The stable characteristic
root, its sign and its modulus should determine whether the leading
correction alternates and which power of $\log s$ replaces
$1/\log s$. Existence of the appropriate positive-measure solutions
must be established before importing the present argument.
\end{enumerate}



% Included source: sections/renewal.tex
\section{The companion sequence: an expansion to every fixed shift}
\label{ren:section}

The comparison of A088713 with A088714 is governed by a large component
in a weighted composition. We prove a general statement because its
hypotheses are exactly those already available for this pair.
Throughout this section, $O_K$ constants may depend on a fixed nonnegative
integer $K$ and on the sequence, but not on $n$.

\begin{theorem}[Renewal expansion for rapidly growing log-convex weights]
\label{ren:theorem}
Suppose $a_0=1$, $a_n>0$, the sequence $(a_n)$ is log-convex, and
$a_n/a_{n-1}\sim n/\log n$. Define the formal series
$A(z)=\sum a_nz^n$, $C(z)=(1-zA(z))^{-1}=\sum c_nz^n$, and
\begin{equation}\label{ren:d}
 d_j=[z^j]C(z)^2.
\end{equation}
For every fixed $K\ge0$,
\begin{equation}\label{ren:expansion}
 c_n=\sum_{j=0}^{K}d_j a_{n-1-j}
       +O_K(a_{n-K-2})\qquad(n\to\infty).
\end{equation}
The error is nonnegative when $n>2K$. This is an asymptotic expansion
in consecutive shifts: $a_{n-j-1}/a_{n-j}\sim\log n/n$ for fixed $j$.
No smooth expansion of the consecutive ratios is assumed.
\end{theorem}

\subsection{A convolution estimate}
Set $w_n=a_{n-1}$ for $n\ge1$, and $q_n=w_{n+1}/w_n$ for $n\ge1$.
Then $w_1=1$, $(q_n)$ is nondecreasing, and $q_n\sim n/\log n$.
For $1\le j\le n-1$ put $T_{n,j}=w_jw_{n-j}$.

\begin{lemma}\label{ren:convolution}
For each fixed integer $J\ge1$,
\begin{equation}\label{ren:tail}
 \sum_{j=J}^{n-J}w_jw_{n-j}=O_J(w_{n-J}).
\end{equation}
There is $\eta>0$ such that
\begin{equation}\label{ren:central}
 \sum_{n/4\le j\le3n/4}w_jw_{n-j}
       =O(e^{-\eta n}w_{n-1}).
\end{equation}
In particular
\begin{equation}\label{ren:smallS}
 S_n:=\frac1{w_n}\sum_{j=1}^{n-1}w_jw_{n-j}
 =O(1/q_{n-1})\longrightarrow0.
\end{equation}
Integer endpoints in these sums are rounded inward.
\end{lemma}
\begin{proof}
Let $L=\lfloor n/4\rfloor$. For $j\le L$,
\[
 \frac{T_{n,j+1}}{T_{n,j}}
   =\frac{q_j}{q_{n-j-1}}
   \le\frac{q_L}{q_{n-L-1}}\longrightarrow\frac13.
\]
Thus the ratio is at most $1/2$, uniformly on this range, for all
sufficiently large $n$. From $L$ to the middle, log-convexity makes
$T_{n,j}$ nonincreasing. The sum from $J$ to $L$ is at most
$2T_{n,J}$; the remainder to the middle is at most
$n2^{-(L-J)}T_{n,J}$. Reflecting across the middle proves
\eqref{ren:tail}, since $w_J$ is fixed. Starting the same bound at $J=1$
shows that the central sum is at most a constant times
$n2^{-L}w_1w_{n-1}$, which implies \eqref{ren:central}.
Taking $J=1$ in \eqref{ren:tail} proves \eqref{ren:smallS}.
\end{proof}

\subsection{Uniform comparison and the unique large part}
The renewal recurrence and composition expansion are
\begin{align}
 c_n&=w_n+\sum_{j=1}^{n-1}w_jc_{n-j},\label{ren:recurrence}\\
 c_n&=\sum_{k=1}^{n}\ \sum_{\substack{s_1+\cdots+s_k=n\\s_i\ge1}}
                   \prod_{i=1}^{k}w_{s_i}.\label{ren:compositions}
\end{align}
Choose $N$ so that $S_n\le1/2$ for $n\ge N$, and choose $M\ge2$
with $c_n\le Mw_n$ for $1\le n<N$. Induction using
\eqref{ren:recurrence} gives
\[
 c_n\le w_n+MS_nw_n\le Mw_n\quad(n\ge N).
\]
Therefore $c_n=O(w_n)$ uniformly. Since
$d_n=2c_n+\sum_{j=1}^{n-1}c_jc_{n-j}$ for $n\ge1$, the same estimate
and \eqref{ren:smallS} give $d_n=O(w_n)$.

A composition has at most one part larger than $n/2$. If that part has
size $n-j$, the parts on its left and right are arbitrary weighted
compositions of combined size $j$. Their generating function is $C(z)^2$.
Consequently the exact contribution of all compositions with such a part is
\begin{equation}\label{ren:giant}
 \sum_{0\le j<n/2}w_{n-j}d_j.
\end{equation}
For $j\ge K+1$, its remaining contribution is bounded by a constant times
\[
 \sum_{j=K+1}^{\lfloor n/2\rfloor}w_jw_{n-j}
       =O_K(w_{n-K-1}).
\]

For a composition with every part at most $n/2$, split it at its first
partial sum at least $n/4$. The partial sum lies in $[n/4,3n/4)$.
The split is uniquely specified by the composition, so the total weight
of this class is at most
\[
 \sum_{n/4\le j\le3n/4}c_jc_{n-j}
 \le M^2\sum_{n/4\le j\le3n/4}w_jw_{n-j}
 =O(e^{-\eta n}w_{n-1}).
\]
For fixed $K$, the quotient $w_{n-1}/w_{n-K-1}$ grows at most
polynomially, as follows from the ratio equivalent. This last contribution
is therefore $o(w_{n-K-1})$. Combining the two classes proves
\eqref{ren:expansion}, with $w_{n-K-1}=a_{n-K-2}$.
Every omitted composition has nonnegative weight, which proves the
sign assertion and completes the proof of Theorem~\ref{ren:theorem}.

\subsection{Consequences for A088713}
Applying Proposition~\ref{gold:foundation} gives a positive answer to
the pinned report's Question 23.6, together with all its fixed-shift
corrections. For this pair,
\[
 (d_0,\ldots,d_6)=(1,2,5,16,64,308,1716).
\]
For example,
\begin{align}
 c_n={}&a_{n-1}+2a_{n-2}+5a_{n-3}+16a_{n-4}
           +O(a_{n-5}),\label{ren:four}\\
 \frac{c_n}{a_{n-1}}
   ={}&1+\frac2{r_{n-1}}
       +\frac5{r_{n-1}r_{n-2}}
       +O\!\left(\frac{(\log n)^3}{n^3}\right).\label{ren:ratio3}
\end{align}
In particular,
\begin{equation}\label{ren:first}
 \frac{c_n}{a_{n-1}}\longrightarrow1,
 \qquad
 \frac{c_n}{a_{n-1}}=1+\frac{2\log n}{n}
                       +o\!\left(\frac{\log n}{n}\right).
\end{equation}
One must retain the actual ratios in \eqref{ren:ratio3} to interpret it
as a second-order statement. Replacing all of them by $n/\log n$
would require more information about their errors than is assumed here.

\begin{corollary}[Additive comparison of shifted consecutive ratios]
\label{ren:additive}
For $r_n=a_n/a_{n-1}$,
\[
 \frac{c_{n+1}}{c_n}
   =r_n+2\left(1-\frac{r_n}{r_{n-1}}\right)+O(1/r_n),
 \qquad
 \frac{c_{n+1}}{c_n}-\frac{a_n}{a_{n-1}}\longrightarrow0.
\]
\end{corollary}
\begin{proof}
Use \eqref{ren:expansion} with $K=1$ at $n$ and $n+1$, divide, and
expand the denominator. Adjacent ratios are asymptotic, so the resulting
quadratic relative errors contribute $O(1/r_n)$ after multiplication
by $r_n$. Finally $r_n/r_{n-1}\to1$.
\end{proof}

The proof uses only ordinary weighted compositions and the leading ratio
law. It does not prove the finer additive expansion of $r_n$ needed for
the Bell-number normalization. The script \texttt{code/verify\_renewal.py}
checks the exact composition identity independently, and records numerical
diagnostics for the consecutive-shift expansion; finite diagnostics are
not used in its proof.


\clearpage
\part{Two-sided fluctuations of unitary-divisor partitions}

% Included source: sections/partitions.tex
\section{Two-sided arithmetic oscillations in A301981 and A301982}
\label{part:sec:unitary}

\subsection{The question and the advance over the repository}

Let
\[
 b_j=\sigma^*(j)=\sum_{\substack{d\mid j\\(d,j/d)=1}}d,
 \qquad
 F_-(z)=\prod_{j\ge1}(1-z^j)^{-b_j},\qquad
 F_+(z)=\prod_{j\ge1}(1+z^j)^{b_j},
\]
and write $a_n^\pm=[z^n]F_\pm(z)$. These are A301981 and A301982,
respectively; the weights form A034448 \cite{oeis301981,oeis301982,oeis034448}.
For the minus sign each color can be repeated; for the plus sign each
individual color can be used once. All coefficients are positive.

Put $\mathcal G$ for the Glaisher--Kinkelin constant and define
\begin{align}
 A_-&=\frac{\pi^2}{6},&
 A_+&=\frac{\pi^2}{8},&
 c_\pm&=3(A_\pm/4)^{1/3},&
 t_{0,\pm}(n)&=(2A_\pm/n)^{1/3},\label{part:eq:scales}\\
 m_n^-&=
 \frac{\mathcal G^6e^{c_-n^{2/3}-1/2}}
 {2\,3^{5/6}\pi^{1/3}n^{5/6}},&
 m_n^+&=
 \frac{e^{c_+n^{2/3}}}
 {2^{4/3}\sqrt3\,\pi^{1/6}n^{2/3}}.\label{part:eq:models}
\end{align}
The OEIS entries label $a_n^\pm\sim m_n^\pm$ as conjectural.
The existing ProveIt report \cite{unitaryprior} already refutes both
equivalents: neither ratio eventually lies in a compact subinterval of
$(0,\infty)$. Its Section~11, question~3, explicitly leaves the direction
of divergence and two-sided oscillation open. Its remark after
Theorem~3.3 stresses that it does not prove both a zero lower limit and
an infinite upper limit.

The result here answers that question. For each sign,
\begin{equation}
 \liminf_{n\to\infty}\frac{a_n^\pm}{m_n^\pm}=0,
 \qquad
 \limsup_{n\to\infty}\frac{a_n^\pm}{m_n^\pm}=+\infty.
 \label{part:eq:ratio-oscillation}
\end{equation}
More quantitatively, there are positive constants such that
$\log(a_n^\pm/m_n^\pm)$ exceeds a positive multiple of $n^{1/4}$
infinitely often and is below a negative multiple of $n^{1/4}$
infinitely often. We give sharper constants depending on an uncancelled
zeta zero and its multiplicity. No Riemann hypothesis, simple-zero
assumption, or expansion over all zeros is used.

The method is the classical positivity argument associated with Landau
\cite{landau1905}; a related Mellin application is in
Diamond--Pintz \cite{diamondpintz2009}. The contribution is its rigorous
application to this explicit repository question, including the
coefficient sampling step and the resulting inverse error. We make no
claim to a new general oscillation principle or an exhaustive historical
priority determination.

\subsection{The prior coefficient bridge, with its analytic justification}

In this subsection fix one sign, and omit it from $A,c,f,R,t_0$ when
convenient. Set
\begin{align}
 f_\pm(t)&=\log F_\pm(e^{-t}),\qquad t>0,\nonumber\\
 d_-&=\tfrac12,&
 B_-&=6\log\mathcal G-\tfrac12-\tfrac12\log(2\pi),&
 d_+&=0,& B_+&=-\tfrac12\log2,\nonumber\\
 R_\pm(t)&=f_\pm(t)-A_\pm t^{-2}-d_\pm\log t-B_\pm.
 \label{part:eq:remainder}
\end{align}
Thus $R_\pm$ is an exactly defined function, not an error symbol with an
assumed sign or bound.

\begin{proposition}[Prior coefficient bridge]\label{part:prop:bridge}
For either sign, as $t\downarrow0$ and $n\to\infty$,
\begin{align}
 (-1)^rf_\pm^{(r)}(t)
 &=A_\pm(r+1)!\,t^{-r-2}+o(t^{-r-1})
 \quad\text{for every fixed }r\ge0,\label{part:eq:derivative}\\
 \log\frac{a_n^\pm}{m_n^\pm}
 &=R_\pm(t_{0,\pm}(n))+o(1).\label{part:eq:bridge}
\end{align}
In particular $h_\pm(x):=R_\pm(1/x)$ satisfies
$h_\pm(x)=o(x)$ and $h_\pm'(x)=o(1)$.
\end{proposition}

These are Lemma~4.1 and Theorem~6.1 of \cite{unitaryprior}.
The following argument records why the bridge is available without an
assumption about zeta zeros inside the critical strip.

\begin{proof}
The elementary bounds
\[
 j\le b_j\le\sigma(j),\qquad
 \sum_{j\le X}b_j\le
 \sum_{d\le X}d\lfloor X/d\rfloor\le X^2
\]
give, by partial summation,
$\sum_j b_jj^re^{-ju}=O_r(u^{-r-2})$ for $u>0$.
Consequently the product logarithms and their derivatives converge
locally uniformly on $\Re t>0$, and
$|f^{(r)}(w)|=O_r(t^{-r-2})$ when $\Re w\ge t/2$.
The classical Dirichlet series \cite[Section~4.9]{toth2017} gives
\begin{align}
 D(s)&=\sum_{j\ge1}b_jj^{-s}
 =\frac{\zeta(s)\zeta(s-1)}{\zeta(2s-1)},\qquad \Re s>2,\nonumber\\
 M_\pm(s)&=\int_0^\infty f_\pm(t)t^{s-1}\,dt
 =\Gamma(s)\chi_\pm(s)
   \frac{\zeta(s+1)\zeta(s)\zeta(s-1)}{\zeta(2s-1)},\label{part:eq:mellin}\\
 \chi_-(s)&=1,\qquad \chi_+(s)=1-2^{-s}.\nonumber
\end{align}
For completeness, at $p^a$ the weight is $1+p^a$, so its Euler factor is
$(1-p^{1-2s})/((1-p^{-s})(1-p^{1-s}))$.
Termwise Mellin integration of
$f_-(t)=\sum_{j,\ell\ge1}b_je^{-j\ell t}/\ell$ proves
\eqref{part:eq:mellin}; $f_+(t)=f_-(t)-f_-(2t)$ proves the plus case.

The kernel for $(-1)^rf^{(r)}(t)$ is obtained by replacing $\Gamma(s)$
by $\Gamma(s+r)$ and $t^{-s}$ by $t^{-s-r}$.
It has residue $A(r+1)!$ at $s=2$.
At $s=1$ the quotient $\zeta(s)/\zeta(2s-1)$ tends to $2$,
so that point is removable. The line $\Re(2s-1)=1$ is zero-free.
The bound $1/\zeta(\sigma+iT)=O(\log(2+|T|))$ for $\sigma\ge1$
and large $|T|$ follows, for example, from Trudgian
\cite{trudgian2015}. Together with the gamma factor and the usual
polynomial bounds for zeta on fixed strips, it permits shifting the
Mellin contour to $\Re s=1$. The remaining integral is
\[
 t^{-r-1}\frac1{2\pi}
 \int_{\mathbb R}K_r(1+iu)e^{iu\log(1/t)}\,du,
 \qquad K_r(1+iu)\in L^1(\mathbb R).
\]
The Riemann--Lebesgue lemma proves \eqref{part:eq:derivative}.

Here are the coefficient details needed below. Let
$\kappa_r(t)=(-1)^rf^{(r)}(t)$, $V(t)=\kappa_2(t)$, and let $t_n$ be
the unique solution of $\kappa_1(t_n)=n$. Existence and uniqueness follow
by realizing each factor as a geometric, or Bernoulli, probability
generating function: $\kappa_1$ is the total mean, $V>0$ is its variance,
and $\kappa_1'=-V$. Equation~\eqref{part:eq:derivative} gives
$t_n\sim t_0$ and $V(t)\sim6At^{-4}$.

The centered characteristic function of this independent sum obeys
\begin{equation}
 |\phi_t(\theta)|
 \le\exp\{-c_1\min(t^{-4}\theta^2,t^{-2})\},
 \qquad |\theta|\le\pi,\label{part:eq:angular}
\end{equation}
with $c_1>0$. Indeed keep the factors with
$N\le j\le2N$, $N=\lfloor1/t\rfloor$.
Their individual logarithmic moduli are at most
$-c_2\sin^2(j\theta/2)$, and $b_j\ge j$.
For $|\theta|\le1/N$, the weighted sine-square sum is
$\gg N^4\theta^2$.
For $1/N\le|\theta|\le K/N$, its quotient by $N^2$ tends uniformly to
the strictly positive integral
$\int_1^2v\sin^2(vN|\theta|/2)\,dv$, for any fixed $K>1$.
For $K/N\le|\theta|\le\pi$, choosing $K>4\pi$ and using
$|\sum_{j=N}^{2N}e^{ij\theta}|\le\pi/|\theta|$ gives
a lower bound $\gg N^2$. This proves \eqref{part:eq:angular}.

On the local scale $z=\theta\sqrt V$, Taylor's theorem gives
$\log\phi_t(z/\sqrt V)=-z^2/2+O(t|z|^3)$ on every fixed $z$ interval.
The bound \eqref{part:eq:angular} supplies an integrable Gaussian
majorant until $|z|\asymp t^{-1}$; beyond it the total integral is
$O(t^{-2}e^{-c_3t^{-2}})$.
Fourier inversion and dominated convergence therefore yield
\begin{equation}
 a_n=\frac{\exp(f(t_n)+nt_n)}{\sqrt{2\pi V(t_n)}}(1+o(1)).
 \label{part:eq:saddle}
\end{equation}
Writing $\kappa_1(t)=2At^{-3}+o(t^{-2})$ and comparing the saddle
equations shows $t_n-t_0=o(t_0^2)$.
The phase $f(t)+nt$ is stationary at $t_n$ and its second derivative
is $O(t_0^{-4})$ between $t_n$ and $t_0$.
Its change on that interval is consequently $o(1)$.
It follows from \eqref{part:eq:saddle} that
\[
 a_n=\frac{t_0^2}{\sqrt{12\pi A}}
       e^{f(t_0)+nt_0}(1+o(1))
     =m_n e^{R(t_0)}(1+o(1)).
\]
The last equality is exact before the factor $1+o(1)$, by
\eqref{part:eq:models} and \eqref{part:eq:remainder}.
Taking logarithms proves \eqref{part:eq:bridge}.
Finally $R(t)=o(t^{-1})$ and $R'(t)=o(t^{-2})$ imply the assertions
about $h$ by the chain rule.
\end{proof}

\subsection{A positivity lemma with the exact pole constant}

The next two lemmas isolate the classical mechanism so that eventual
one-sided bounds are not confused with absolute bounds.

\begin{lemma}[Landau's positivity principle]\label{part:lem:landau}
Let $q:[1,\infty)\to[0,\infty)$ be locally integrable and of at most
polynomial growth. If the real abscissa of convergence $\alpha$ of
$Q(s)=\int_1^\infty q(x)x^{-s-1}\,dx$ is finite, then $Q$ has a
singularity at the real point $s=\alpha$.
\end{lemma}

\begin{proof}
Suppose $Q$ were holomorphic near $\alpha$. Set $x=e^u$.
For $s_1>\alpha$,
\[
 (-1)^jQ^{(j)}(s_1)
 =\int_0^\infty u^j q(e^u)e^{-s_1u}\,du\ge0.
\]
Holomorphy in the right half-plane together with a disc about $\alpha$
allows $s_1$ sufficiently close to $\alpha$ and $\delta>s_1-\alpha$
such that the Taylor series at $s_1$ converges at $s_1-\delta$.
For example, if the disc about $\alpha$ has radius $r$, take
$s_1=\alpha+r/4$ and $r/4<\delta<r/2$.
Tonelli's theorem then identifies the finite value of that Taylor series with
\[
 \sum_{j\ge0}\frac{\delta^j}{j!}
 \int_0^\infty u^jq(e^u)e^{-s_1u}\,du
 =\int_0^\infty q(e^u)e^{-(s_1-\delta)u}\,du.
\]
This contradicts the definition of $\alpha$.
\end{proof}

\begin{lemma}[A quantitative two-sided consequence]
\label{part:lem:quantitative}
Let $h:[1,\infty)\to\mathbb R$ be locally integrable and of at most
polynomial growth. Suppose
$H(s)=\int_1^\infty h(x)x^{-s-1}\,dx$, initially defined far to the
right, continues meromorphically to $\Re s>0$ and is holomorphic near
every positive real point. Suppose $s_*=\sigma_*+i\tau_*$, where
$\sigma_*>0$ and $\tau_*\ne0$, is a pole of order $m\ge1$, with
\[
 C_*=\lim_{s\to s_*}(s-s_*)^mH(s)\ne0.
\]
Then
\begin{align}
 \limsup_{x\to\infty}
 \frac{h(x)}{x^{\sigma_*}(\log x)^{m-1}}
 &\ge\frac{|C_*|}{(m-1)!},\nonumber\\
 \liminf_{x\to\infty}
 \frac{h(x)}{x^{\sigma_*}(\log x)^{m-1}}
 &\le-\frac{|C_*|}{(m-1)!}.
 \label{part:eq:quantitative-landau}
\end{align}
\end{lemma}

\begin{proof}
If the first inequality failed, choose
$0<c<|C_*|/(m-1)!$ exceeding the alleged upper limit.
Then $q(x)=cx^{\sigma_*}(\log x)^{m-1}-h(x)\ge0$ for $x\ge X$,
with $X>1$. Set $q=0$ on $[1,X)$. Its Mellin transform has the
meromorphic continuation
\begin{equation}
 Q(s)=\frac{c(m-1)!}{(s-\sigma_*)^m}-H(s)+E_X(s),
 \label{part:eq:envelope-transform}
\end{equation}
where $E_X$ is entire, being the integral over a finite interval.
This expression is holomorphic near every real $s>\sigma_*$.
Lemma~\ref{part:lem:landau} shows that the abscissa of convergence of
the nonnegative integral is at most $\sigma_*$; if that abscissa is
$-\infty$, the conclusion is immediate.
Hence for real $\sigma>\sigma_*$,
$|Q(\sigma+i\tau_*)|\le Q(\sigma)$.
The identity \eqref{part:eq:envelope-transform} holds throughout this
half-plane by uniqueness of continuation.
Multiply that inequality by $(\sigma-\sigma_*)^m$ and let
$\sigma\downarrow\sigma_*$. The limits of its two sides are,
respectively, $|C_*|$ and $c(m-1)!$:
the artificial pole in \eqref{part:eq:envelope-transform} is regular
at $s_*$, whereas $H$ is regular at the real point $\sigma_*$.
The resulting inequality contradicts the choice of $c$.
Apply the same argument to $-h$ to obtain the lower limit.
\end{proof}

No rightmost-pole hypothesis appears here. If a further pole lies to
the right, the assumed one-sided envelope already contradicts the
holomorphy of its nonnegative integral in that larger half-plane.
The proof needs only a single nonreal pole whose leading Laurent
coefficient is nonzero.

\subsection{Noncancellation and the coefficient theorem}

Choose a nontrivial zeta zero $\rho=\beta+i\gamma$ of least positive
imaginary part, choosing its reflection about the critical line so
that $\beta\ge1/2$. Such a least height exists because zeta zeros are
isolated, bounded-height portions of the closed critical strip are
compact away from the pole at $1$, and there are no real nontrivial
zeros. Let its multiplicity be $m$ and put
\begin{equation}
 s_*=\frac{1+\rho}{2}=\sigma_*+i\tau_*,
 \qquad \frac34\le\sigma_*<1.
 \label{part:eq:selected-zero}
\end{equation}
Only the usual location, symmetry, and existence facts about zeta
zeros are needed \cite{dlmfzeros}.

\begin{lemma}[The surviving pole and its constant]\label{part:lem:pole}
For each sign, $M_\pm$ in \eqref{part:eq:mellin} has a pole of order
$m$ at $s_*$, and its leading Laurent coefficient is
\begin{equation}
 C_\pm=
 \frac{m!\,\Gamma(s_*)\chi_\pm(s_*)
       \zeta(s_*+1)\zeta(s_*)\zeta(s_*-1)}
      {2^m\zeta^{(m)}(\rho)}
 \ne0.
 \label{part:eq:laurent}
\end{equation}
\end{lemma}

\begin{proof}
The denominator has leading term
$2^m\zeta^{(m)}(\rho)(s-s_*)^m/m!$.
Gamma has neither a zero nor a pole at $s_*$, and
$\zeta(s_*+1)\ne0$ because $\Re(s_*+1)>1$.
If $\zeta(s_*)=0$, its positive imaginary part $\gamma/2$ would
contradict the selection of $\gamma$.
The point $s_*-1$ has real part in $[-1/4,0)$ and nonzero imaginary
part, so it is neither a trivial nor a nontrivial zero of zeta.
Finally $|2^{-s_*}|<1$, so $\chi_+(s_*)\ne0$.
Division of the leading terms proves the formula.
\end{proof}

\begin{theorem}[Quantified two-sided oscillation]\label{part:thm:main}
With \eqref{part:eq:selected-zero}--\eqref{part:eq:laurent}, define
\[
 K_\pm=
 \frac{|C_\pm|}{(m-1)!}\,
 (2A_\pm)^{-\sigma_*/3}\,3^{\,1-m}>0,
 \qquad
 W(n)=n^{\sigma_*/3}(\log n)^{m-1}.
\]
For each sign separately,
\begin{align}
 \limsup_{n\to\infty}
 \frac{\log(a_n^\pm/m_n^\pm)}{W(n)}&\ge K_\pm,\nonumber\\
 \liminf_{n\to\infty}
 \frac{\log(a_n^\pm/m_n^\pm)}{W(n)}&\le-K_\pm.
 \label{part:eq:main}
\end{align}
The same inequalities hold after replacing
$\log(a_n^\pm/m_n^\pm)$ by $\log a_n^\pm-c_\pm n^{2/3}$.
Consequently \eqref{part:eq:ratio-oscillation} holds, and both centered
logarithms are $\Omega_\pm(n^{1/4})$.
\end{theorem}

\begin{proof}
For $h_\pm(x)=R_\pm(1/x)$ and $\Re s>2$, splitting the Mellin integral
at $1$ gives
\begin{equation}
 H_\pm(s)=M_\pm(s)-\frac{A_\pm}{s-2}
               +\frac{d_\pm}{s^2}-\frac{B_\pm}{s}-T_\pm(s),
 \quad
 T_\pm(s)=\int_1^\infty f_\pm(t)t^{s-1}\,dt .
 \label{part:eq:H}
\end{equation}
The function $T_\pm$ is entire because $f_\pm(t)$ decays
exponentially at infinity. Formula~\eqref{part:eq:H} is therefore a
meromorphic continuation to $\Re s>0$ and has the same pole and
Laurent coefficient as $M_\pm$ at $s_*$.

It is holomorphic near every positive real point.
The pole at $2$ has been removed, the point $1$ is removable as
already checked, and the rational terms have poles only at $0$.
At every other positive real $s$, $\zeta(2s-1)$ has no zero:
its argument is in $(-1,\infty)$, an interval without real zeta
zeros. The pole of that denominator at $s=1$ is already covered.
There are no other positive real singularities of the numerator.
Proposition~\ref{part:prop:bridge} supplies the polynomial growth of
$h_\pm$. Thus Lemma~\ref{part:lem:quantitative} applies.

It remains to preserve both signs when sampling.
Let $x_n=(n/(2A_\pm))^{1/3}$.
Consecutive values differ by $O(x_n^{-2})$.
For any large $x$, choose $n$ with $x_n\le x<x_{n+1}$.
The derivative estimate $h_\pm'=o(1)$ gives
$h_\pm(x)-h_\pm(x_n)=o(x_n^{-2})$.
Also $x/x_n\to1$ and $\log x/\log x_n\to1$.
The continuous upper and lower bounds therefore hold along these
discrete samples. Since
\[
 x_n^{\sigma_*}(\log x_n)^{m-1}
 =(2A_\pm)^{-\sigma_*/3}3^{\,1-m}
 n^{\sigma_*/3}(\log n)^{m-1}(1+o(1)),
\]
\eqref{part:eq:bridge} proves \eqref{part:eq:main}.
The difference between the two centered logarithms is $O(\log n)$,
which is negligible on the displayed scale.
Finally $\sigma_*/3\ge1/4$ and $m\ge1$; each of the positive and
negative subsequences is unbounded in magnitude.
Exponentiation proves \eqref{part:eq:ratio-oscillation}.
\end{proof}

The constants in \eqref{part:eq:main} are rigorous lower bounds on
oscillation amplitudes, not exact amplitudes or estimates of the
first index where a sign occurs. In particular the proof does not
assert that a single zero dominates the full remainder.

\subsection{A consequence for inversion}

Extend either expression $m_n^\pm$ in \eqref{part:eq:models} to a
positive function $m^\pm(x)$ for real $x>0$, and let $X_\pm(y)$ be
its inverse on its eventually increasing tail.
Let
\[
 N_\pm(y)=\min\{n\ge0:a_n^\pm\ge y\}.
\]
The earlier report proves that $a_n^\pm$ is eventually strictly
increasing \cite[Theorem~8.2]{unitaryprior}. This is the only
additional prior input in the following corollary.

\begin{corollary}[Oscillating error of the pure inverse]
\label{part:cor:inverse}
Put
\[
 L_\pm=
 \frac{2^{\,1-m}|C_\pm|}
 {(m-1)!\,(3A_\pm)^{(\sigma_*+1)/2}}>0.
\]
Then
\begin{align}
 \limsup_{y\to\infty}
 \frac{N_\pm(y)-X_\pm(y)}
 {(\log y)^{(\sigma_*+1)/2}(\log\log y)^{m-1}}
 &\ge L_\pm,\nonumber\\
 \liminf_{y\to\infty}
 \frac{N_\pm(y)-X_\pm(y)}
 {(\log y)^{(\sigma_*+1)/2}(\log\log y)^{m-1}}
 &\le-L_\pm.
 \label{part:eq:inverse}
\end{align}
In particular this inverse error is
$\Omega_\pm((\log y)^{7/8})$.
\end{corollary}

\begin{proof}
It suffices to use $y_n=a_n^\pm$, for which $N_\pm(y_n)=n$ once
the increasing tail exceeds the finite initial prefix.
Write $g(x)=\log m^\pm(x)$ and $E_n=\log(a_n^\pm/m_n^\pm)$.
The mean value theorem gives
\[
 X_\pm(y_n)-n=\frac{E_n}{g'(\xi_n)}.
\]
Here $\log a_n^\pm=c_\pm n^{2/3}+o(n^{1/3})$ by
Proposition~\ref{part:prop:bridge}, so $X_\pm(y_n)/n\to1$,
$\xi_n/n\to1$, and
$g'(\xi_n)=(2A_\pm)^{1/3}n^{-1/3}(1+o(1))$.
Apply Theorem~\ref{part:thm:main}, with the sign reversed.
Since $n\sim(\log y_n/c_\pm)^{3/2}$, conversion of the scale yields
the stated constant:
\[
 K_\pm(2A_\pm)^{-1/3}
 c_\pm^{-(\sigma_*+1)/2}(3/2)^{m-1}
 =L_\pm.
\]
\end{proof}

Thus a leading logarithmic equivalent can support a correct relative
inverse scale while its pure-model inverse has large errors of both
signs. The exact-product inverse already constructed in
\cite{unitaryprior} retains the arithmetic remainder and is unaffected
by this obstruction.

\subsection{Scope and further questions}

The proven advance is the resolution of the repository's two-sided
oscillation question, with the stronger $n^{1/4}$ lower scale and an
inverse consequence. The existing refutation, Mellin kernel,
least-height noncancellation argument, and coefficient bridge are
credited prior inputs; the bridge and noncancellation have been
rederived above to expose every analytic hypothesis.
The accompanying \texttt{code/verify\_partitions.py} compares direct
unitary-divisor enumeration with prime-power factorization, then
compares binomial product expansion with the logarithmic-derivative
recurrence. All comparisons pass through $n=160$, including the first
21 displayed OEIS terms of each sequence. The recorded exact output is
\texttt{data/partition\_verification.json}. These computations check
definitions and finite data; they cannot certify either infinitely
occurring sign.

Several substantive questions remain.
\begin{enumerate}
\item Determine how often each sign occurs. The proof produces
arbitrarily large indices of both signs; it supplies no natural or
logarithmic density and no bound on gaps between sign changes.
\item Obtain effective indices and certified positive lower bounds
for $K_\pm$. This requires certified information on a selected zero
and quantitative control of the positivity argument, beyond simply
evaluating \eqref{part:eq:laurent} in floating arithmetic.
\item Determine the sharp oscillation scale and amplitude.
The bound uses one surviving pole, which need not be rightmost.
Even under RH, a limiting distribution would require control of the
combined zero contributions and possible relations between their
imaginary parts.
\item Construct a rigorously truncated zero-residue formula with an
explicit remainder. Pole locations alone do not justify summing
all residues, exchanging limits, or claiming a convergent
oscillatory expansion.
\item Study joint behavior of the two signs. Their exact relation
$f_+(t)=f_-(t)-f_-(2t)$ links the remainders at different scales,
but the separate $\liminf/\limsup$ assertions above imply no
simultaneous sign pattern.
\end{enumerate}


\clearpage

% Included source: sections/verification_and_agenda.tex
\section{Verification, reproducibility, and interpretation}
\label{verify:section}

The accompanying archive contains the complete \LaTeX{} source, the PDF,
four verification programs, their exact or interval output, and a short
reproduction guide. The programs do not access the network. The stored
OEIS prefixes and the repository pin make their finite reference data
explicit. Standard-library integer and rational arithmetic suffice for
the renewal and partition checks. The dynamics certificate additionally
uses correctly rounded decimal elementary functions with outward
enclosure. The map program uses SymPy for exact symbolic identities and
mpmath only for numerical asymptotic diagnostics.

\subsection{What was checked}

\begin{center}
\small
\begin{tabularx}{\textwidth}{@{}p{0.26\textwidth}X@{}}
\toprule
Program & Verification performed\\
\midrule
\texttt{verify\_maps.py}
& Checks both positive coefficient formulas against all 20 displayed
OEIS terms; compares the formulas through $n=100$; independently
enumerates connected rotation systems through four edges; verifies the
bridge substitutions, elimination relation, and Puiseux coefficients
symbolically.\\[4pt]
\texttt{verify\_dynamics.py}
& Generates $a_0,\ldots,a_{31}$ exactly, encloses $F(1/100)$ by its
finite Stieltjes remainders, propagates 115 inverse steps with outward
rounding, and encloses both invariant-series tails.\\[4pt]
\texttt{verify\_renewal.py}
& Compares two independent formal coefficient constructions through
degree 25; checks the stored OEIS prefixes; enumerates all 16,383
compositions of sizes 1 through 14 and checks the large-part identity;
computes diagnostics through $n=260$.\\[4pt]
\texttt{verify\_partitions.py}
& Computes unitary-divisor weights by divisor enumeration and by
prime-power factorization; compares direct product expansion with
the logarithmic-derivative recurrence through $n=160$, including all
21 stored initial terms of each sequence.\\
\bottomrule
\end{tabularx}
\end{center}

Separate mathematical reviews checked the root normalization and
genus derivative in Part I; the global-inversion hypotheses, boundary
limits, and zero descent in the density theorem; the invariant
coordinates, analytic reversion, and interval enclosure in the
golden-ratio expansion; the convolution bounds in the renewal theorem;
and the noncancellation, positivity argument, saddle bridge, and
inverse scaling in Part III. These reviews are internal checks of a
new manuscript. They are not external refereeing or machine-checked
formal proofs.

\subsection{A certificate with a mathematical role}

The interval calculation in the dynamics section has a stronger role
than a plot or an agreement of finite prefixes. The proof reduces
nonvanishing of the analytic profile to the strict inequality
$d(x)>|D(x)-d(x)|$ at one actual point of the OEIS solution.
The enclosure establishes
\[
 0.00003580<d(x)<0.00003582,\qquad
 |D(x)-d(x)|<7.9\cdot10^{-8}.
\]
It also encloses the corresponding profile value by
\[
 0.00095032738531560
 <P\!\left(\log_\varphi\Psi(x)\right)
 <0.00095032738531562.
\]
Both endpoints are deliberately rounded outwards from the more precise
enclosures in \path{data/dynamics_certificate.json}.
The coefficient truncation error, elementary-function rounding, orbit
propagation, and invariant-series tails are all included. The narrow
interval stored for the explicit error \emph{bound} does not claim to
enclose the actual absolute error.

This proves $P\not\equiv0$. Antiperiodicity then proves the occurrence
of both signs at arbitrarily large arguments. It does not locate the
maximum of $P$, determine its least positive period, or establish
how many zeros it has in a period.

\subsection{Diagnostics that do not enter the proofs}

The renewal expansion implies that, for fixed $K$,
\[
 E_{K,n}:=
 \frac{c_n-\sum_{j=0}^{K}d_j a_{n-1-j}}{a_{n-K-2}}
 \longrightarrow d_{K+1}.
\]
This follows by applying Theorem~\ref{ren:theorem} with $K+1$.
The exact-integer calculations give the following rounded values.
\begin{center}
\small
\begin{tabular}{@{}r r r r r r@{}}
\toprule
$n$ & $c_n/a_{n-1}$ & $E_{0,n}$ & $E_{1,n}$ & $E_{2,n}$ & $E_{3,n}$\\
\midrule
20  &1.235499&2.692517&7.716257&29.465201&142.036296\\
50  &1.118554&2.315888&6.089803&20.729599&88.744437\\
100 &1.069135&2.178820&5.594920&18.478299&76.410619\\
200 &1.039769&2.101282&5.330955&17.352619&70.645714\\
260 &1.032154&2.081579&5.265425&17.079934&69.282715\\
\midrule
Limit &1&2&5&16&64\\
\bottomrule
\end{tabular}
\end{center}
The slow approach explains why finite ratios alone would be a poor
substitute for the large-part proof. The maps likewise have a visible
$n^{-1/2}$ correction: $(27/256)^n b_n$ is approximately $0.03376725$,
$0.03771030$, and $0.03968765$ at $n=32,128,512$, respectively, while its
limit is $1/24\approx0.04166667$. The three-term approximation has the
proved error $O(n^{-5/2})$; the numerical agreement is only a check of
the implemented constants.

There is deliberately no plot presented as evidence of the partition
sign theorem. The argument is an infinite-subsequence result. A short
coefficient prefix need not exhibit its full oscillation scale, and
an uncertified finite sum over zeta zeros would add no proof of the
claimed lower and upper limits.

\section{A proposed research program}
\label{future:section}

The following directions are ordered by how directly they continue the
proved structures. Each states a concrete objective and the obstacle
that remains. The local questions at the end of Parts I and III supply
additional variants.

\subsection{Finish the finer Bell comparison}

The leading coefficient ratio and the new fixed-shift renewal expansion
leave the principal coefficient-normalization problem open. In the
notation of the pinned report, let $\mathsf B_n$ be the Bell number and
let $W(n)>0$ satisfy $W(n)e^{W(n)}=n$. The conjecture is that a finite
constant $C>0$ exists with
\begin{equation}\label{future:bell}
 a_n\sim C\,\mathsf B_n
          \exp\!\bigl(W(n)^2+3W(n)\bigr).
\end{equation}
A proved sufficient condition from \cite{goldenprior} is, for some
fixed $p\ge0$,
\begin{equation}\label{future:ratio}
 \frac{a_n}{a_{n-1}}
 =\frac n{W(n)}+2+\frac{W(n)}{2(1+W(n))^2}
             +O\!\left(\frac{W(n)^p}{n}\right).
\end{equation}
After division by the leading ratio, the indicated error is summable;
this is the precision needed to produce a nonzero multiplicative
constant.

The new endpoint analysis concerns the measures near zero. High
moments instead probe their remote tails, so it does not directly
imply \eqref{future:ratio}. A useful next target is uniform control of
the moment ratios on windows $n-j$ with $j=O(W(n))$, together with
the complementary part of the coefficient recurrence. Part III of
the pinned report formulates an exact scalar defect whose bound
$O(W(n)^q/n^2)$ would suffice. Its definition and precision should
be retained when attempting this step. Proving only an $o(1)$ error
in a formal local approximation will not settle
\eqref{future:bell}.

\subsection{Quantify the analytic oscillation profile}

The construction supplies an analytic, canonical, nonzero function
$P$, with a rigorous evaluator along inverse orbits. It leaves several
well-posed numerical and analytic problems:
\begin{enumerate}
\item Produce a certified enclosure of $P$ and $P'$ on an entire
period, and use it to determine the number and multiplicity of its
zeros and extrema.
\item Bound the width of a complex strip of analyticity explicitly.
This would turn existence of exponentially decreasing Fourier
coefficients into effective truncation bounds.
\item Determine whether the first allowed Fourier harmonic dominates.
Antiperiodicity allows only the frequencies
$\exp(\pi i(2k+1)t)$; it does not by itself prove such dominance.
\item Bound an explicit threshold for convergence of the inverse-log
series and describe the exponentially smaller correction terms.
The present series is convergent for sufficiently large $\log s$,
but it is not an exact representation of the transform without its
stated exponentially small remainder.
\end{enumerate}
These are refinements of a proved oscillation, not a renewed question
about whether oscillation exists.

\subsection{From positive analytic densities to effective tail laws}

The density theorem removes possible gaps, atoms, and internal zeros,
and gives precise behavior at the origin. The large-$x$ behavior of
$p(x)$ and $q(x)$ remains the natural bridge to the finer moment
asymptotics. One can ask for a tail equivalent with an error uniform
under the exponential tilts that localize moments near their saddle.
Such a result would supply information that positivity and real
analyticity alone cannot provide.

A more immediate computational objective is a certified solver for the
boundary relation
\[
 \int\frac{t^2\,d\nu(t)}{(u-t)^2+v(u)^2}=1
\]
on compact positive intervals, coupled to the density formulas of
Section~\ref{reg:subsec:analyticity}. Finite moments and transform
identities should be used to bound the truncated tails. An algorithm
with an explicit enclosure for the discarded measure would allow
quantitative plots and tests of global monotonicity or global
log-convexity. The derivative equivalents already give
$(\log p)''(x)\sim\beta/x^2$ and
$(\log q)''(x)\sim\alpha/x^2$ at zero, so both densities are strictly
log-convex sufficiently near zero. The two global shape properties
remain unresolved.

\subsection{Find a direct map bijection and extend the genus calculation}

The exact square has the integral factor
$\frac14\binom{4n}{n}=\binom{4n-1}{n-1}$.
This suggests a direct decomposition of a bridgeless rooted toroidal
map into a pair of independently described objects of total size $n$.
A bijection should explain the root and the positive convolution,
rather than merely reproducing the same algebraic generating function.
It would also clarify whether a face or vertex statistic has a natural
refinement on the two factors.

The multigenus bridge equation already contains a separate algebraic
research program. Insert known unrestricted genus-$g$ map series,
extract the coefficient of the genus marker, and retain every
derivative caused by the substituted argument. Determine rational
forms in the common planar parameter and systematic singular
exponents for the bridgeless series. Higher-genus expressions are
not obtained by a naive substitution into the unrestricted series;
Part I displays the missing contribution already at genus one.

\subsection{Make the partition oscillations effective}

The two-sided theorem provides explicit symbolic lower constants
from a genuine Mellin pole, but no first index at which either bound
must occur. A useful next result would combine a certified zero
enclosure, a residue enclosure, and a quantitative version of the
positivity argument to produce a finite search interval containing
an index of each sign.

A different question concerns statistics rather than existence:
do the normalized logarithmic errors have limiting distributions,
and what are the densities or gaps of sign changes? Any answer must
control the combined zero contributions and the remainder of a
truncated explicit formula. The exact identity
$f_+(t)=f_-(t)-f_-(2t)$ also makes joint sign patterns a distinct
problem. Separate oscillation of the two sequences gives no theorem
about their simultaneous behavior.

\subsection{Broaden the two general mechanisms}

Two proofs invite abstraction beyond these particular OEIS entries.
First, the renewal theorem uses a uniform contraction of
$q_j/q_{n-j-1}$ away from the endpoints. One can seek sharp hypotheses
on positive log-convex weights, including ratio laws regularly varying
with an arbitrary positive index, under which the same full
fixed-shift expansion holds. One should also identify what replaces
the exponential central bound at the boundary of that regime.

Second, inverse functional equations can produce a recurrence with
one expanding eigenvalue and one negative contracting eigenvalue.
The present Fibonacci recurrence turns this spectral structure into
an analytic antiperiodic correction. A general theorem should state
which perturbations preserve analytic conjugacies, what the resonant
corrections are when additional eigenvalues are present, and how to
certify that a particular combinatorial solution has a nonzero
stable coordinate.

\subsection{Formalization and an independent submission path}

The exact bridge decomposition, the positive coefficient formulas,
and the renewal convolution estimates are comparatively self-contained
targets for formalization. The interval certificate can be ported
to a small formally checked rational-enclosure kernel. The density
and Mellin arguments require a broader analytic foundation, making
them later targets.

For external mathematical review, the three Parts can be separated
into articles with their dependencies preserved. An OEIS update for
A343093 should cite the complete rooted-map proof and its normalization;
updates for the partition entries should distinguish the earlier
refutation from the new two-sided strengthening. No update or
submission has been made as part of this work. A bibliographic
comparison with the older map-enumeration literature and independent
refereeing remain appropriate before asserting historical priority.


\clearpage
\begingroup
\raggedright
\phantomsection
\addcontentsline{toc}{section}{References}
\begin{thebibliography}{99}

% Included source: bibliography_items.tex
\bibitem{map:oeis343093}
The OEIS Foundation Inc., \emph{A343093: Number of rooted toroidal maps
with $n$ edges and no isthmuses}, entry by Andrew Howroyd; conjectural
generating function by Peter Bala, July 24, 2025.
\url{https://oeis.org/A343093}. Accessed October 4, 2026 (UTC).

\bibitem{map:bcr1988}
E.~A. Bender, E.~R. Canfield, and R.~W. Robinson,
\emph{The enumeration of maps on the torus and the projective plane},
Canadian Mathematical Bulletin \textbf{31} (1988), 257--271.
\url{https://doi.org/10.4153/CMB-1988-039-4}.
Exact inputs: equations (4.20) and (4.22), p.~268.

\bibitem{map:cyz2019}
J. Courtiel, K. Yeats, and N. Zeilberger,
\emph{Connected chord diagrams and bridgeless maps},
Electronic Journal of Combinatorics \textbf{26}(4) (2019), Paper P4.37,
56 pp. \url{https://doi.org/10.37236/7400};
\url{https://arxiv.org/abs/1611.04611}.
See Proposition~23, pp.~17--18, with the auxiliary dangling-root
normalization explained in the present text.

\bibitem{goldenprior}
ProveIt research-report collection,
\emph{Bell-Scale Growth of OEIS A088714}, merged article, Parts I--III,
September 20--October 2, 2026. AI-assisted unrefereed research report.
\href{https://github.com/VladimirReshetnikov/ProveIt/blob/eaf08931cbfd0132dd36ab82a7743ccb02506dc7/SetTheory/Cardinals/docs/reports/generating-functions-and-asymptotics/oeis-sequence-asymptotics/a088714-bell-scale-growth/article.tex}{Pinned source in Vladimir Reshetnikov's ProveIt repository}.
Snapshot \texttt{eaf08931cbfd0132dd36ab82a7743ccb02506dc7}.
The analytic foundations and Questions 23.1, 23.3, 23.6 are in Part II;
Part III identifies the remaining fine Bell-comparison estimates.

\bibitem{oeisA088714}
The OEIS Foundation Inc., \emph{A088714}.
\url{https://oeis.org/A088714}. Accessed October 4, 2026 (UTC).

\bibitem{oeisA088713}
The OEIS Foundation Inc., \emph{A088713}.
\url{https://oeis.org/A088713}. Accessed October 4, 2026 (UTC).

\bibitem{HuangWang2022}
H.-W. Huang and J.-C. Wang,
\emph{Regularity results for free L\'evy processes},
Advances in Mathematics \textbf{402} (2022), Article 108323.
\url{https://doi.org/10.1016/j.aim.2022.108323};
\url{https://arxiv.org/abs/2203.00421}.
The global inversion and boundary results used here are
Propositions 3.1--3.2.

\bibitem{python-decimal}
Python Software Foundation,
\emph{Decimal fixed-point and floating-point arithmetic},
Python 3 library documentation.
\url{https://docs.python.org/3/library/decimal.html}.
Correct rounding of \texttt{exp}, \texttt{ln}, and \texttt{sqrt}, together
with adjacent representable values, is used for outward enclosure.
Accessed October 4, 2026 (UTC).

\bibitem{unitaryprior}
ProveIt research-report collection,
\emph{Unitary divisor sum colored partitions and their asymptotics},
October 2, 2026. AI-assisted unrefereed research report.
\href{https://github.com/VladimirReshetnikov/ProveIt/blob/eaf08931cbfd0132dd36ab82a7743ccb02506dc7/SetTheory/Cardinals/docs/reports/generating-functions-and-asymptotics/oeis-sequence-asymptotics/a301981-unitary-divisor-partitions/article.tex}{Pinned source in Vladimir Reshetnikov's ProveIt repository}.
Snapshot \texttt{eaf08931cbfd0132dd36ab82a7743ccb02506dc7}.

\bibitem{oeis301981}
The OEIS Foundation Inc.,
\emph{A301981: Euler transform of A034448},
entry by V. Kot\v{e}\v{s}ovec, March 30, 2018.
\url{https://oeis.org/A301981}. Accessed October 4, 2026 (UTC).

\bibitem{oeis301982}
The OEIS Foundation Inc.,
\emph{A301982: Expansion of
$\prod_{k\ge1}(1+x^k)^{\mathrm{A034448}(k)}$},
entry by V. Kot\v{e}\v{s}ovec, March 30, 2018.
\url{https://oeis.org/A301982}. Accessed October 4, 2026 (UTC).

\bibitem{oeis034448}
The OEIS Foundation Inc.,
\emph{A034448: Sum of unitary divisors}.
\url{https://oeis.org/A034448}.

\bibitem{landau1905}
E. Landau, \"Uber einen Satz von Tschebyschef,
\emph{Mathematische Annalen} \textbf{61} (1905), 527--550.
\url{https://geodesic.mathdoc.fr/item/MAN_1905__61_158244/}.

\bibitem{diamondpintz2009}
H. G. Diamond and J. Pintz, Oscillation of Mertens' product formula,
\emph{Journal de th\'eorie des nombres de Bordeaux}
\textbf{21} (2009), no.~3, 523--533.
\url{https://doi.org/10.5802/jtnb.687}.

\bibitem{toth2017}
L. T\'oth, Alternating sums concerning multiplicative arithmetic functions,
\emph{Journal of Integer Sequences} \textbf{20} (2017), Article 17.2.1.
Section~4.9, p.~28, records the ordinary unitary-divisor Dirichlet series.
\url{https://cs.uwaterloo.ca/journals/JIS/VOL20/Toth/toth25.pdf}.

\bibitem{trudgian2015}
T. S. Trudgian, Explicit bounds on the logarithmic derivative and
the reciprocal of the Riemann zeta-function,
\emph{Functiones et Approximatio Commentarii Mathematici}
\textbf{52} (2015), no.~2, 253--261.
\url{https://doi.org/10.7169/facm/2015.52.2.5}.

\bibitem{dlmfzeros}
NIST Digital Library of Mathematical Functions,
Section~25.10, \emph{Zeros of the Riemann zeta function}.
\url{https://dlmf.nist.gov/25.10}.


\end{thebibliography}
\endgroup
\end{document}
